% concatenated from Belavkin_POC_restructured.tex

\documentclass[10pt,a4paper,reqno]{amsart}

\usepackage[utf8]{inputenc}
\usepackage[T1]{fontenc}
\usepackage{amsthm,amsmath,amsfonts,amssymb,mathtools}
\usepackage[numbers]{natbib}
\usepackage{hyperref}
\hypersetup{hidelinks}
\usepackage{enumitem}
\usepackage{microtype}
\usepackage{mathrsfs}
\usepackage[a4paper,margin=1in]{geometry}

\allowdisplaybreaks
\numberwithin{equation}{section}

\theoremstyle{plain}
\newtheorem{theorem}{Theorem}[section]
\newtheorem{proposition}[theorem]{Proposition}
\newtheorem{lemma}[theorem]{Lemma}
\newtheorem{corollary}[theorem]{Corollary}
\newtheorem{assumption}[theorem]{Assumption}
\newtheorem{hypothesis}[theorem]{Hypothesis}

\theoremstyle{definition}
\newtheorem{definition}[theorem]{Definition}

\theoremstyle{remark}
\newtheorem{remark}[theorem]{Remark}

\renewcommand{\d}{\mathrm d}
\renewcommand{\i}{\mathrm i}
\newcommand{\Tr}{\mathrm{Tr}}
\newcommand{\Law}{\mathrm{Law}}
\newcommand{\Ent}{\mathrm{Ent}}
\newcommand{\TV}{\mathrm{TV}}
\newcommand{\E}{\mathbb E}
\newcommand{\Pp}{\mathbb P}
\newcommand{\Q}{\mathbb Q}
\newcommand{\cS}{\mathcal S}
\newcommand{\cL}{\mathcal L}
\newcommand{\cF}{\mathcal F}
\newcommand{\cD}{\mathcal D}
\newcommand{\cC}{\mathcal C}
\newcommand{\cE}{\mathcal E}
\newcommand{\cY}{\mathcal Y}
\newcommand{\HH}{\mathcal H}
\newcommand{\KK}{\mathcal K}
\newcommand{\1}{\mathbf 1}

\title[Belavkin equations beyond pure states]
{Propagation of Chaos for Belavkin Equations Beyond Pure States}

\author[G. Guo]{Gaoyue Guo}
\address{Universit\'e Paris-Saclay, CentraleSup\'elec, MICS and CNRS FR-3487}
\email{gaoyue.guo@centralesupelec.fr}

\subjclass[2020]{Primary 60K35, 60H10; Secondary 60F05, 81Q93, 81S25, 93E11}
\keywords{Belavkin equation, quantum filtering, 
propagation of chaos, McKean--Vlasov diffusion, stochastic BBGKY hierarchy,
Zakai equation, relative entropy}

\begin{document}



\begin{abstract}
We use probabilistic and stochastic-analysis methods to prove trace-norm
propagation of chaos for finite-dimensional quantum mean-field systems governed
by Belavkin equations.  The particles are density matrices,
interact through a mean-field Hamiltonian, and are continuously monitored
through independent diffusive observation channels.  The limiting dynamics is a
nonlinear matrix-valued McKean--Vlasov diffusion, random through its local
observation record and coupled through the deterministic averaged state.

The main result treats arbitrary one-particle density matrices, including mixed
states, and both perfect and inefficient measurement regimes.  Under strong
tensorization of the initial data, every fixed marginal converges uniformly on
compact time intervals to the tensor product of the nonlinear limiting filters,
with an explicit quantitative bound.  The proof combines purification, fully
observed dilation, conditional expectation, relative entropy, and uniform
stability of the associated Zakai equations.  In the skew-adjoint measurement
case, exterior observation noises disappear from the marginal equations and a
stochastic BBGKY hierarchy is recovered.  Under only marginal chaoticity of
permutation-invariant initial states, we prove convergence of fixed marginals by
an iteration of this hierarchy.
\end{abstract}

\maketitle

\setcounter{tocdepth}{2}\tableofcontents





\section{Introduction}
\label{sec:introduction}

Propagation of chaos is a law of large numbers for interacting particle
systems.  It asserts that, in the large-particle limit, every fixed
subcollection of particles becomes asymptotically independent and is described
by copies of an effective nonlinear one-particle dynamics.  In the present
paper the particles are quantum states.  More precisely, they are density
matrices on a finite-dimensional Hilbert space, interact through a mean-field
Hamiltonian, and are continuously monitored through independent diffusive
observation channels.  The limiting equation is a nonlinear matrix-valued
McKean--Vlasov diffusion.

The problem lies at the interface between quantum mechanics, stochastic
analysis and mean-field probability.  On the quantum side, the state of one
particle is not a point in a Euclidean phase space, but a positive trace-one
operator, which may be viewed as a noncommutative analogue of a probability
distribution.  The interaction is generated by commutators with Hamiltonians,
and the irreversible part of the dynamics is governed by Lindblad operators.
On the probabilistic side, the limiting object is a nonlinear stochastic
process: each particle is driven by its own observation record, while the
interaction with the population enters through the averaged one-particle
state.  Thus the quantum structure is encoded in the coefficients, whereas the
large-population limit has the form of a McKean--Vlasov law of large numbers.

Mean-field limits for closed quantum systems have been studied extensively.
Among the available methods are BBGKY hierarchies, counting functionals,
Heisenberg expansions, Egorov-type arguments, quantum empirical measures and
relative entropy.  We refer, among others, to Spohn \cite{spohn1980},
Bardos--Golse--Mauser \cite{bardos2000},
Bardos--Erd\H{o}s--Golse--Mauser--Yau \cite{bardos2002},
Fr\"ohlich--Graffi--Schwarz \cite{frohlich2007}, Ammari--Nier
\cite{Ammari2009}, Knowles--Pickl \cite{knowlespickl2010}, Pickl
\cite{pickl2011}, Golse--Paul \cite{golsepaul2017,golsepaul2019}, and
Guo--Liang--Wang \cite{Guo2026}.  In the closed quantum setting, the
\(N\)-particle density matrix solves the deterministic von Neumann equation
\[
  \i\partial_t\rho_t^N=[H^N,\rho_t^N].
\]
This equation is linear at the \(N\)-particle level, and after taking partial
traces the reduced density matrices satisfy a deterministic BBGKY hierarchy.
The equation for the \(n\)-particle marginal is coupled to the
\((n+1)\)-particle marginal only through the interaction term.  This triangular
structure is one of the basic mechanisms behind the derivation of Hartree
equations.

A similar simplification remains available for open quantum systems when the
environment is not continuously observed.  In that case the density matrix
satisfies a deterministic Lindblad equation.  The local dissipative
terms are trace-preserving on the variables which are traced out, and hence do
not create exterior stochastic records in the marginal equations.  The problem
is still quantum and may be technically difficult, but the reduced dynamics is
again deterministic.

The situation changes substantially for continuously observed systems.  We
study Belavkin equations, also called stochastic master equations, which
describe the conditional state of a quantum system given a measurement record,  
see, e.g.  Bouten--van Handel--James \cite{bouten2007}, Pellegrini
\cite{pellegrini2008}, and Mirrahimi--van Handel
\cite{mirrahimivanhandel2007}.  The normalized equation is nonlinear and
stochastic.  Each particle carries its own observation record, and the
stochastic coefficient depends on the full conditional state.  Consequently,
after one traces out particles, the Brownian motions attached to the traced
particles do not disappear in general.  The marginal equation for the observed
subsystem still contains information coming from exterior observation
channels.  This phenomenon, a kind of stochastic coupling created by conditioning on the
observation records, is the basic obstruction to a direct adaptation of the
BBGKY method.



The same obstruction appears in other forms.  In the Heisenberg picture,
testing against a local observable produces covariance terms with remote
observed coordinates, and the evolution of these covariances involves higher
centered correlations.  The standard empirical-measure strategy from
classical mean-field theory is also not directly available.  Indeed, for
classical particles one can encode the configuration by the random empirical
measure,  prove tightness of these empirical measures, and identify the limit by a
martingale problem.  For a general quantum state on \(\HH^{\otimes N}\),
there is no literal random empirical measure of the individual particles:
the state is a noncommutative object, and its marginals are obtained by
partial traces rather than by projecting a random configuration.  Quantum
empirical-measure methods provide a powerful replacement in closed systems,
but in the present continuously observed setting they do not close the
stochastic hierarchy at a fixed order.  Passing to the unnormalized Zakai
equation removes the explicit normalization, but the exterior observation
records still carry information about the particles that have been traced out.
Entropy methods are also delicate, since natural quantum relative-entropy
estimates may require faithfulness or uniform absolute-continuity assumptions
which are not stable in pure-state regimes.  Thus the main difficulty is not
only the mean-field interaction, but the interaction between marginalization
and stochastic conditioning.


This is one reason why probability and stochastic analysis enter the present
work in an essential way.  The problem is motivated by quantum filtering and
quantum control, but the mechanism of the proof is probabilistic: conditioning,
Girsanov transforms, relative entropy, conditional factorization, and stability
of Zakai equations are used to recover a law of large numbers in a
noncommutative state space.  In this sense, the paper uses stochastic analysis
to address a question arising naturally from continuously monitored and
controlled quantum systems.

Besides its intrinsic probabilistic interest, propagation of chaos also has a
model-reduction meaning for quantum simulation.  If \(d=\dim(\HH)\), then an
exact simulation of the \(N\)-particle conditional state requires matrices on
\(\HH^{\otimes N}\), hence of size \(d^N\times d^N\).  Even if one is
interested only in a fixed marginal \(\rho_t^{N:n}\), the Belavkin dynamics
does not provide a closed equation for that marginal in general.  Direct
simulation therefore retains an exponential dependence on \(N\).  A
propagation-of-chaos result justifies replacing \(\rho_t^{N:n}\) by
\[
  \Gamma_t^n
  =
  \gamma_t^1\otimes\cdots\otimes\gamma_t^n,
\]
where the processes \(\gamma^j\) are independent copies of an effective
one-particle nonlinear filter.  This replaces the many-body conditional
dynamics by matrix-valued McKean--Vlasov diffusions, together with the
deterministic evolution of the averaged state.  Such an approximation is
naturally accessible to Monte Carlo simulation and avoids the full many-body
Hilbert-space dimension.

This numerical viewpoint is consistent with the role of quantum trajectories
as stochastic unravellings of deterministic quantum dynamical semigroups, a
method widely used for Monte Carlo evaluation of open quantum dynamics.  It is
also connected with quantum feedback control, quantum dynamic games and
quantum mean-field games, where the state used for feedback is itself a
conditional state produced by observation.  For a recent systematic account of
quantum filtering, propagation of chaos and these applications, we refer to the survey  \cite{kolokoltsov2026survey} of
Kolokoltsov.  From this perspective, propagation
of chaos for Belavkin equations is a dynamic law of large numbers for
continuously monitored quantum systems, and it provides a probabilistic
foundation for replacing large observed or controlled quantum systems by
effective nonlinear filtering equations.

A first propagation-of-chaos theory for continuously observed quantum
many-particle systems was developed by Kolokoltsov
\cite{kolokoltsov2021,kolokoltsov2022}.  These works treat the pure-state,
perfect-efficiency regime.  In that case the evolution can be represented at
the wave-function level and controlled by a counting functional.  This method
uses the fact that the trace distance to a rank-one state is controlled by a
scalar quantity measuring the number of particles outside a prescribed
one-particle wave function.  When the initial state is mixed, or when the
measurement efficiency is strictly less than one, the conditional state is no
longer expected to remain pure.  The wave-function representation and the
associated counting functional are then no longer available.

The main purpose of the present paper is to prove propagation of chaos beyond
this pure-state setting.  Our first result treats arbitrary mixed one-particle
initial states and every measurement efficiency \(0<\eta\le1\), in finite
dimension.  If the initial \(N\)-particle state is close in trace norm to a
tensor product,
\[
  \delta_N
  :=
  \left\|
    \rho_0^N-\gamma_0^{\otimes N}
  \right\|_1
  \longrightarrow0,
\]
then every fixed marginal converges, uniformly on compact time intervals,
toward the corresponding tensor product of nonlinear limiting particles.  The
estimate is quantitative and contains both the initial tensorization error and
an explicit rate in the particle number.

The proof combines three probabilistic reductions.  First, at perfect
efficiency, mixed one-particle states are handled by purification.  One
realizes a mixed density matrix as the partial trace of a pure state on an
enlarged Hilbert space, lifts the Hamiltonian, interaction and measurement
operators, applies the pure-state propagation result, and finally traces out
the auxiliary variables.  The point is that this quantum-information
construction is compatible with both the interacting Belavkin dynamics and the
mean-field limit.

Second, inefficient observations are treated by a fully observed dilation.  The
unobserved part of each output channel is represented by an additional
Brownian motion, so that the enlarged system has perfect total observation.
Conditioning on the actually observed coordinates recovers the inefficient
equation.  At finite \(N\), however, this conditioning does not preserve a
product structure.  We compare the law of the interacting observed paths with
the product one-particle law by a compensating Girsanov transform.  The
resulting relative-entropy estimate, together with the entropy chain rule and
Pinsker's inequality, yields approximate conditional factorization.

Third, the exact tensor-product assumption on the full initial state is removed
through a stability estimate for the associated linear Zakai equation.  A
direct comparison of two nonlinear \(N\)-particle Belavkin equations would
lead to constants growing with the number of particles, because the stochastic
coefficient contains \(mN\) noise terms.  Instead, we use a common linear
reference equation.  Its solution map is positive and preserves the trace in
expectation.  These two properties give an \(N\)-uniform trace-norm stability
estimate for the unnormalized equations, and the normalization is handled by a
change of measure and a total-variation estimate.

Our second result concerns an exceptional skew-adjoint regime.  When
\[
  (L^{(k)})^*=-L^{(k)},
  \qquad k=1,\ldots,m,
\]
the innovation coefficient becomes linear.  In this case the exterior Brownian
motions vanish from the marginal equations by cyclicity of the partial trace,
and one recovers a stochastic BBGKY hierarchy.  This allows us to assume only
marginal chaoticity of the initial states, namely convergence of every fixed
initial marginal to the corresponding tensor product.  Under this weaker
assumption we prove convergence of every fixed marginal, without an explicit
rate, by closing the stochastic hierarchy through an iteration argument in the
spirit of Bardos--Golse--Mauser \cite{bardos2000}.

The paper is organized as follows.  Section~\ref{sec:framework} introduces the
finite-dimensional framework, states the main results, and gives a detailed
comparison with closed quantum systems and with Kolokoltsov's convergence
result.  Section~\ref{sec:purification} proves the perfect-efficiency
product-state estimate by purification.  Section~\ref{sec:inefficient} treats
inefficient observation through fully observed dilation, conditional
expectation and relative entropy.  Section~\ref{sec:initial-perturbation}
proves the \(N\)-uniform Zakai stability estimate and transfers the
product-state bounds to approximately tensorized initial data.  Finally,
Section~\ref{sec:skew-bbgky} proves the weak-tensorization result in the
skew-adjoint case by a stochastic BBGKY hierarchy.







\bigskip



\paragraph{\bf Acknowledgments.} This work was supported by ANR-25-CE40-0714 (MATH-SPA).  The author thanks Hao Liang and Zhenfu Wang for helpful discussions.  




\section{Mathematical framework and main results}
\label{sec:framework}


Throughout the paper, \(\HH\) is reserved for a finite-dimensional complex Hilbert
space.  It represents the state space of a single quantum system,
for instance a spin, a qubit, a qudit, or a finite-level atom.  We keep this
notation abstract and do not choose coordinates.  The relevant state variable is
not only a \emph{wave function}  taking values in $\HH$,  but more generally identified by a \emph{density operator} on \(\HH\).
This is essential for continuously monitored open quantum systems.  Indeed, an
indirect measurement is usually modeled by coupling the system to an auxiliary
probe or environment and then measuring the outgoing probe.  Conditioning on
the random measurement record produces a stochastic evolution of the system
state, often called a quantum trajectory,  see e.g. Pellegrini \cite{pellegrini2008} and Bouten--van Handel--James
\cite{bouten2007}.

For \(N\) identical particles,  the one-particle space \(\HH\) is replaced by their
tensor product
\(
 \HH^{\otimes N}.
\)
One-particle observables are lifted to individual tensor factors, two-particle
interactions are lifted to pairs of tensor factors, and reduced states of
subsystems are obtained by partial traces.  These elementary operations provide
the notation used to formulate the \(N\)-particle Belavkin equation and its
mean-field limit introduced later. 

%%%%%%%%%%%%%%%%%%%%%%%%%%%%%%%%%%%%%%%%%%%%%%%%%%%%%%%%%%%%%%%%%%%%%%%%%%%%%%%
\subsection{Quantum states and observables}
\label{subsec:states-and-composite-systems}
%%%%%%%%%%%%%%%%%%%%%%%%%%%%%%%%%%%%%%%%%%%%%%%%%%%%%%%%%%%%%%%%%%%%%%%%%%%%%%%

Denote by $\cL(\HH)$ the space of linear operators on $\HH$,  and for
$X\in\cL(\HH)$, let
\[
  \|X\|
  :=
  \sup_{\psi\in\HH:\,\|\psi\|_{\HH}\le1}
  \|X\psi\|_{\HH},
  \qquad
  \|X\|_1
  :=
  \Tr\bigl(\sqrt{X^*X}\bigr),
  \qquad
  \|X\|_2
  :=
  \sqrt{\Tr(X^*X)}
\]
denote, respectively, the \emph{operator},  \emph{trace} and \emph{Hilbert--Schmidt norms}.  Since
$\HH$ is finite-dimensional, these norms are equivalent. Namely,  
\[
  \|X\|
  \le
  \|X\|_2
  \le
  \|X\|_1
  \le
  \dim(\HH)\,\|X\|,\qquad \mbox{ for all } X\in\cL(\HH).
\]
The corresponding \emph{Hilbert--Schmidt inner product} is given as 
\(
  \langle X,Y\rangle_{\mathrm{HS}}
  :=
  \Tr(X^*Y). 
\)
For $X,Y\in\cL(\HH)$, we use the notation
\[
  [X,Y]:=XY-YX,
  \qquad
  \{X,Y\}:=XY+YX
\]
for the \emph{commutator} and \emph{anticommutator}.  Further,  we introduce the following subsets of $\cL(\HH)$:
\begin{align*}
  \cL^*(\HH)
  &:=
  \{X\in\cL(\HH):X^*=X\},
  \\
  \cL^+(\HH)
  &:=
  \{X\in\cL^*(\HH):X\ge0\},
  \\
  \cS(\HH)
  &:=
  \{\rho\in\cL^+(\HH):\Tr(\rho)=1\}.
\end{align*}
Elements of $\cS(\HH)$ are called density operators,  or density matrices in the finite-dimensional setting,  and represent \emph{quantum states},  while elements of $\cL^*(\HH)$ stand for  \emph{observables}.  If
$A\in\cL^*(\HH)$ and $\rho\in\cS(\HH)$, then
\(
  \Tr(\rho A)
\)
stands for the expectation of the observable $A$ in the state $\rho$.

\medskip

A state $\rho\in\cS(\HH)$ is called \emph{pure} if it has rank one,  i.e.,  $\operatorname{rank}(\rho)=1$.  
Equivalently, there exists some unit vector $\psi\in\HH$,  identified as a wave function,  such that
\(
  \rho=|\psi\rangle\langle\psi|,
\)
where $ |\psi\rangle\langle\psi| \in \cL(\HH)$ denotes the \emph{projection operator} given as
\[
  |\psi\rangle\langle\psi|\,x
  :=
  \langle\psi,x\rangle_{\HH}\,\psi,
  \qquad
  x\in\HH.
\]
Thus a wave function determines a pure density operator, and two wave functions
which differ only by a global phase determine the same physical state.  A state
which is not pure is called \emph{mixed}.  Every mixed state $\rho$ admits a spectral
representation
\[
  \rho
  =
  \sum_{r=1}^{\operatorname{rank}(\rho)}
  \lambda_r|\psi_r\rangle\langle\psi_r|,
  \qquad
  \lambda_r>0,
  \qquad
  \sum_{r=1}^{\operatorname{rank}(\rho)}\lambda_r=1.
\]
The numbers
\(\lambda_r\) are scalar weights, whereas the operator \(\rho\) is the
state itself.  A mixed state may describe classical statistical uncertainty,
but it may also arise as the reduced state of a larger entangled quantum
system.  This distinction is one reason why density operators, rather than
only wave functions, are the natural state variables for open quantum
systems.

\medskip

For $N\ge1$, let
\(
  [N]:=\{1,\ldots,N\}
\) 
and
  \(
  \HH_N:=\HH^{\otimes N}
\)
be the $N$-tensor product of $\HH$.  The space \(\HH_N\) is the Hilbert space of \(N\) labeled copies of the
one-particle system.  Product density operators describe uncorrelated
configurations, while a general element of $\cS(\HH_N)$ may contain both
classical correlations and quantum entanglement.  

\medskip

If \(J=\{j_1<\cdots<j_r\}\subset[N]\), we set
\[
  \HH_J
  :=
  \HH_{j_1}\otimes\cdots\otimes\HH_{j_r},
  \qquad
  \HH_{J^c}
  :=
  \bigotimes_{j\in J^c} \HH_j,
\]
where \(\HH_j\) denotes the \(j\)-th copy of \(\HH\). We use
the convention that the tensor product over the empty set is \(\mathbb C\).
Thus \(\HH_J\) is canonically identified with \(\HH_{|J|}\).  Let 
\[
  \Tr_{J^c}:
  \cL(\HH_N)
  \longrightarrow
  \cL(\HH_{J})
\]
denote the \emph{partial trace}  which keeps the tensor factors indexed by $J$ and
traces out the complementary factors.  It is characterized by the duality
relation
\[
  \Tr\left(
    A\,\Tr_{J^c}(X)
  \right)
  =
  \Tr\left(
    (A\otimes I_{\HH_{J^c}})X
  \right)
\]
for every local observable $A$ on the factors indexed by $J$, after the
canonical reshuffling of tensor factors.  When no confusion is possible, we identify \(\HH_J\) with
\(\HH_{|J|}\).  In particular, the partial trace
extracts the state seen by a subsystem.

\medskip

For $A\in\cL(\HH)$ and $j\in[N]$, let $A_j\in\cL(\HH_N)$ denote the lift
of $A$ to the $j$-th tensor factor:
\[
  A_j
  :=
  I_{\HH}^{\otimes(j-1)}
  \otimes A\otimes
  I_{\HH}^{\otimes(N-j)}.
\]
Similarly, if $O \in\cL(\HH\otimes\HH)$ and $1\le j<k\le N$, then
$O_{jk}$ denotes the lift of $O$ to the $j$-th and $k$-th tensor
factors.



\subsection{Quantum system under continuous indirect measurements}
\label{subsec:continuously-monitored-systems}





We fix a self-adjoint one-particle Hamiltonian
\[
  h=h^*\in\cL^*(\HH),
\]
a finite family of measurement,  or coupling,  operators
\[
  L^{(1)},\ldots,L^{(m)}\in\cL(\HH),
  \qquad
  m\ge1,
\]
and a self-adjoint two-body  interaction
\[
  V=V^*\in\cL^*(\HH\otimes\HH)
\]
which is symmetric,  i.e.,   invariant under exchange of the two tensor factors.  The operator
\(h\) generates the coherent evolution of one isolated particle.  The operator
\(V\) describes the interaction between two particles.  The operator
\(L^{(k)}\) describes the coupling of a particle to the \(k\)-th output
channel,  and the associated observed quadrature is
\(L^{(k)}+(L^{(k)})^*\).  Thus the same operator \(L^{(k)}\) leads to  both
the dissipative back-action of the environment and the stochastic update
generated by the measurement record.

\medskip

For $\sigma\in\cS(\HH)$, define the effective one-particle interaction
\begin{equation}
  V^\sigma
  :=
  \Tr_{\{2\}}
  \bigl(
    V(I_{\HH}\otimes\sigma)
  \bigr),
  \label{eq:effective-potential}
\end{equation}
where we recall that $\Tr_{\{2\}}$ denotes the partial trace over the second tensor factor. The operator \(V^\sigma\) is the mean field created by infinitely many particles of the same  state 
\(\sigma\).  
The cyclicity of the partial trace on the traced factor gives
\begin{equation}
  \Tr_{\{2\}}
  \bigl(
    [V,\gamma\otimes\sigma]
  \bigr)
  =
  [V^\sigma,\gamma],
  \qquad \mbox{ for all }
  \gamma,\sigma\in\cS(\HH).
  \label{eq:effective-commutator}
\end{equation}
For every $L\in\cL(\HH)$, define the maps $\cD_L$,  $\cC_L$ and $\cE_L$ on $\cL(\HH)$ by
\begin{align}
  \cD_L(X)
  &:=
  LXL^*-\frac12\{L^*L,X\},
  \label{eq:D-def}
  \\
  \cC_L(X)
  &:=
  LX+XL^*,
  \label{eq:C-def}
  \\
  \cE_L(X)
  &:=
  \cC_L(X)
  -
  \Tr\bigl((L+L^*)X\bigr)X.
  \label{eq:E-def}
\end{align}
The map $\cD_L$ is the Lindblad dissipator generated by the coupling \(L\).
It describes the irreversible back-action of the environment on the state.  The
map $\cC_L$ is the linear coefficient appearing in the (unnormalized) Zakai 
equation, while $\cE_L$ is its normalized correction
which updates the conditional state after a measurement,  and
keeps the trace equal to one.  We use the same notation with a subscript \(j\)
when the corresponding operator acts on the \(j\)-th tensor factor.

\medskip

The mean-field $N$-particle Hamiltonian is
\begin{equation}
  H^N
  :=
  \sum_{j=1}^N h_j
  +
  \frac1N
  \sum_{1\le j<\ell\le N}
  V_{j\ell}.
  \label{eq:HN}
\end{equation}
The factor $N^{-1}$ is the mean-field scaling: every particle interacts weakly
with every other particle, but the total interaction experienced by one
particle remains of order one as $N\longrightarrow\infty$.

\medskip

We now describe the observation model.  The system is coupled to an
auxiliary probe or output field, and the outgoing probe is measured.  The
conditional state, given the observation record, then follows a stochastic
trajectory.  In the diffusive case such trajectories are described by Belavkin
equations, or equivalently by stochastic master equations,  see, e.g.  Bouten--van Handel--James \cite{bouten2007} and Pellegrini
\cite{pellegrini2008}.  We assume that each particle is monitored
separately through the same \(m\) output channels.

\medskip

Let \(\eta\in(0,1]\) be the common\footnote{For notational simplicity we take a common efficiency \(\eta\).  The same
arguments extend to channel-dependent efficiencies
\(\eta_1,\ldots,\eta_m\in(0,1]\), with constants depending on these parameters.} measurement efficiency.   Perfect detection
corresponds to \(\eta=1\), whereas \(0<\eta<1\) means that only part of the
information carried by the output fields is observed. 

\medskip 

For an initial state $\rho_0^N\in\cS(\HH_N)$, the state of the \(N\)-particle system evolves according to 
the Belavkin equation
\begin{equation}
\begin{aligned}
  \d\rho_t^N
  &=
  -\i[H^N,\rho_t^N]\,\d t
  +
  \sum_{k=1}^m\sum_{j=1}^N
  \cD_{L_j^{(k)}}(\rho_t^N)\,\d t
 +
  \sqrt{\eta}
  \sum_{k=1}^m\sum_{j=1}^N
  \cE_{L_j^{(k)}}(\rho_t^N)\,\d W_t^{k,j},
\end{aligned}
  \label{eq:N-Belavkin}
\end{equation}
where $W^{k,j}$ are independent real-valued Brownian motions.  The Hamiltonian commutator gives the coherent many-body evolution.  The
Lindblad terms describe the full back-action of the couplings to the output
fields, whether or not those fields are detected.  The stochastic terms encode
the information gained from the observed part of the outputs,  corresponding to the fact that 
$\eta$ appears in the innovation coefficients.  

\medskip

There is an important pure-state reduction.  If \(\eta=1\) and the initial
state $\rho^N_0$ is pure, then the solution of \eqref{eq:N-Belavkin} remains pure,  saying  \(\rho_t^N=|\Psi_t^N\rangle\langle\Psi_t^N|\),  where the wave function
\(\Psi_t^N\) satisfies the \emph{nonlinear stochastic Schr\"odinger equation}.
This formulation with respect to the  \(N\)-particle wave function  has been extensively studied, both in finite
dimension and in infinite-dimensional Hilbert spaces,  see e.g. Mora--Rebolledo \cite{mora2008}.  Under the pure-state assumption, recent work of de Bouard--Guo--H\'erouard
\cite{bouard2026} also derives infinite-dimensional mean-field limits in this
wave-function setting.  For  the general Belavkin equation \eqref{eq:N-Belavkin},  finite-dimensional well-posedness with respect to the strong solution  was proved by Pellegrini \cite{pellegrini2008}  for
\(\dim(\HH)=2\) and by Mirrahimi--van Handel
\cite{mirrahimivanhandel2007}  for general
finite-dimensional \(\HH\).  More recently,  Kolokoltsov obtained
infinite-dimensional well-posedness results in
\citep{kolokoltsov2025a,kolokoltsov2025b}.




\medskip

For $N$ interacting particles, however, this conditional state belongs to
$\cS(\HH_N)$, whose dimension grows exponentially in $N$.  A
propagation-of-chaos theorem is therefore a model-reduction result which justifies
replacing the many-particle conditional dynamics, at the level of any fixed
number of particles, by independent copies of an effective one-particle filter.  As $N\longrightarrow \infty$, we recover formally the limiting dynamics,  i.e.,  matrix-valued McKean--Vlasov equation
\begin{equation}
\begin{aligned}
  \d\gamma_t^j
  =
  \left(
    -\i[h+V^{\xi_t},\gamma_t^j]
    +
    \sum_{k=1}^m
    \cD_{L^{(k)}}(\gamma_t^j)
  \right)\d t+
  \sqrt{\eta}
  \sum_{k=1}^m
  \cE_{L^{(k)}}(\gamma_t^j)\,\d W_t^{k,j},
  \qquad
  \xi_t:=\E[\gamma_t^j],
\end{aligned}
  \label{eq:MF-Belavkin}
\end{equation}
where $ \gamma_0^j=
  \gamma_0\in\cS(\HH)$.  
The processes $\gamma^1,\gamma^2,\ldots$ are independent and identically distributed.  Their interaction
occurs only through the deterministic mean state $\xi_t$.  Taking expectations
in \eqref{eq:MF-Belavkin} gives the Hartree equation 
\begin{equation}
  \frac{\d}{\d t}\xi_t
  =
  -\i[h+V^{\xi_t},\xi_t]
  +
  \sum_{k=1}^m
  \cD_{L^{(k)}}(\xi_t),
  \qquad
  \xi_0=\gamma_0.
  \label{eq:mean-state-equation}
\end{equation}
Thus the random limiting state is driven by the local measurement record,  whereas 
its mean-field Hamiltonian is determined by the average state of the ensemble.  Well-posedness of   \eqref{eq:MF-Belavkin} is derived by 
Amini--Amini--Chalal--Guo
 \cite{Amini2025} for finite-dimensional $\HH$,  while it remains  open  for general infinite-dimensional $\HH$.  In the present finite-dimensional setting,  we shall use this well-posedness without further comment.

\medskip 

For $J\subset [N]$,  define the $J$-marginal by
\[
  \rho_t^{N:J}
  :=
  \Tr_{J^c}(\rho_t^N)
\]
and write in particular 
\(
  \rho_t^{N:n}
  :=
  \Tr_{[N]\setminus[n]}(\rho_t^N)
\) for 
$1\le n\le N$.   The product of the limiting one-particle processes is given as 
\(
  \Gamma_t^n
  :=
  \gamma_t^1\otimes\cdots\otimes\gamma_t^n.
\)

\begin{definition}[Propagation of chaos]
\label{def:propagation-of-chaos}
Let $\rho_0^N\in\cS(\HH_N)$ and $\gamma_0\in\cS(\HH)$.  Propagation of chaos is said to hold if, for every
fixed $n\ge1$ and $T>0$,
\begin{equation}
  \sup_{0\le t\le T}
  \E\left[
    \left\|
      \rho_t^{N:n}-\Gamma_t^n
    \right\|_1
  \right]
  \longrightarrow0
  \qquad
  \text{as }N\longrightarrow \infty.
  \label{eq:propagation-of-chaos}
\end{equation}
\end{definition}

The convergence in \eqref{eq:propagation-of-chaos} says that every fixed
collection of particles is asymptotically described by independent copies of the
nonlinear one-particle filter.  Under the pathwise well-posedness of \eqref{eq:N-Belavkin} and  \eqref{eq:MF-Belavkin},  the
joint law of the synchronously coupled processes is uniquely determined.  Hence
the expectation in \eqref{eq:propagation-of-chaos} is independent of the
particular filtered probability space on which the coupled equations  \eqref{eq:N-Belavkin},  \eqref{eq:MF-Belavkin} are
defined.

\medskip

For later use, let $\mathfrak S_N$ denote the permutation group of $[N]$.
For $\pi\in\mathfrak S_N$, let $U_\pi\in\cL(\HH_N)$ be the unitary operator 
defined by
\[
  U_\pi(x_1\otimes\cdots\otimes x_N)
  :=
  x_{\pi^{-1}(1)}\otimes\cdots\otimes x_{\pi^{-1}(N)}.
\]
\begin{definition}[Permutation invariance]
\label{def:exchangeability}
A  state $\rho^N\in\cS(\HH_N)$ is called permutation invariant if 
\[
  U_\pi\rho^N U_\pi^*=\rho^N,
  \qquad
  \mbox{ for all } \pi\in\mathfrak S_N.
\]
This is the labelled-particle analogue of the symmetry appearing in bosonic
closed quantum systems, although we do not restrict the state space to the symmetric
tensor product.
\end{definition}


This symmetry is preserved by the Belavkin dynamics in the sense of distributions.  Indeed, if
$\rho_0^N$ permutation invariant, then
$U_\pi\rho_t^N U_\pi^*$ solves the same equation as $\rho_t^N$, but driven by
the relabelled Brownian family
\[
  (W^{k,\pi^{-1}(j)})_{1\le k\le m,\ 1\le j\le N}.
\]
Since this relabelled family has the same law as the original Brownian family,
and since $H^N$ and the measurement coefficients are invariant under particle
relabeling, uniqueness in law implies
\[
  (U_\pi\rho_t^N U_\pi^*)_{t\ge0}
  \stackrel{\mathrm{law}}{=}
  (\rho_t^N)_{t\ge0}.
\]
We shall refer to this property as \emph{exchangeability}.  The same terminology will
be used for coupled families, such as
\[
  \bigl(R^N,Z^1,\ldots,Z^N\bigr),
\]
where exchangeability means invariance in law under the simultaneous relabeling
\[
  \bigl(R^N,Z^1,\ldots,Z^N\bigr)
  \longmapsto
  \bigl(U_\pi R^N U_\pi^*,Z^{\pi^{-1}(1)},\ldots,Z^{\pi^{-1}(N)}\bigr).
\]
In particular, exact tensor-product initial data
$\rho_0^N=\gamma_0^{\otimes N}$ yields the exchangeability properties used
below.








\subsection{Main results}
\label{subsec:main-results}
%%%%%%%%%%%%%%%%%%%%%%%%%%%%%%%%%%%%%%%%%%%%%%%%%%%%%%%%%%%%%%%%%%%%%%%%%%%%%%%

We now state the main results.  

\begin{theorem}
\label{thm:main-approx}
Assume that
\begin{equation}
  \delta_N
  :=
  \left\|
    \rho_0^N-\gamma_0^{\otimes N}
  \right\|_1
  \longrightarrow0.
  \label{eq:strong-tensorization}
\end{equation}
Then propagation of chaos holds.  More precisely, for every fixed $n\ge1$ and
$T>0$, there exists a constant $C>0$ depending on $n, T,  h,V,L^{(1)},\ldots,L^{(m)},\eta $ 
such that
\begin{equation}
\begin{aligned}
  \sup_{0\le t\le T}
  \E\left[
    \left\|
      \rho_t^{N:n}
      -
      \Gamma_t^n
    \right\|_1
  \right]
   \le C\delta_N+ 
     \frac{C}{N^{1/4}}{\mathbf 1}_{\{\eta=1\}} + \frac{C}{N^{1/16}}{\mathbf 1}_{\{0<\eta<1\}}.
\end{aligned}
  \label{eq:main-approx-bound}
\end{equation}
\end{theorem}

\begin{remark}
Theorem~\ref{thm:main-approx}  simultaneously removes the two restrictions in Kolokoltsov 
\citep{kolokoltsov2021,kolokoltsov2022}: the initial one-particle state may
be genuinely mixed, and the measurement efficiency $\eta$ may take values in 
$(0,1)$.  In addition,  no exact permutation invariance of $\rho_0^N$ is assumed.
\end{remark}


\begin{theorem}
\label{thm:main-anti}
Assume that $L^{(1)},\ldots, L^{(m)}$ are skew-adjoint and 
 the initial states $\rho^N_0$ are permutation invariant.  Then the propagation of chaos holds provided that,  for every fixed $n\ge 1$, 
\begin{equation}
  \delta_{N}^n
  :=
  \left\|
    \rho_0^{N:n}-\gamma_0^{\otimes n}
  \right\|_1
  \longrightarrow 0.
  \label{eq:weak-tensorization}
\end{equation}
\end{theorem}

\begin{remark}
In the skew-adjoint case, the innovation coefficients $\cE_{L_j^{(k)}}$ become linear and the
exterior noises disappear from the marginal equations.  A stochastic BBGKY
argument can therefore propagate the weaker marginal chaoticity assumption
\eqref{eq:weak-tensorization}.
\end{remark}



\subsection{Comparison with closed systems and Kolokoltsov's convergence result}
\label{subsec:difficulty-and-open-problem}


When all measurement operators vanish, equation \eqref{eq:N-Belavkin}
reduces to the deterministic von Neumann equation for a closed quantum system:
\[
  \i\partial_t\rho_t^N=[H^N,\rho_t^N].
\]
In this closed setting, and in particular for infinite-dimensional
Schr\"odinger dynamics with \(\HH=L^2(\mathbb R^d;\mathbb C)\), propagation
of chaos and the derivation of Hartree equations have been established by
several complementary methods: deviation estimates,  e.g. 
Knowles--Pickl \citep{knowlespickl2010} and Pickl \citep{pickl2011}; BBGKY
hierarchies,  e.g.  Spohn \citep{spohn1980} and
Bardos--Golse--Mauser \citep{bardos2000}; Egorov-type theorems,  e.g.  
Fr\"ohlich--Graffi--Schwarz \citep{frohlich2007}; expansions of the
Heisenberg dynamics,  e.g.   Ammari--Nier \citep{Ammari2009}; quantum
empirical-measure techniques,  e.g.  
Golse--Paul \citep{golsepaul2017,golsepaul2019}; and quantum relative
entropy,  e.g.  Guo--Liang--Wang \citep{Guo2026}.  Although these
approaches are quite different, they exploit two structural features of closed
systems.  First, the \(N\)-particle evolution is linear and unitary.  Second,
after taking a partial trace, the deterministic BBGKY hierarchy is triangular:
the equation for an \(n\)-particle marginal is coupled to the
\((n+1)\)-particle marginal only through the interaction term.

Continuous indirect measurements remove both simplifying features.
\begin{itemize}
\item The Belavkin equation is nonlinear and stochastic, and its stochastic
integrals are attached to the individual observation records.  In particular,
after tracing out particles, the Brownian motions corresponding to the traced
particles do not disappear.  This is the main reason why the closed-system
methods do not extend by a routine perturbation argument.  The BBGKY hierarchy
is no longer a single deterministic triangular hierarchy: its coefficients
involve labeled higher-order marginals and exterior observation channels.

\item In the Heisenberg picture, testing against a local observable transforms
the same terms into covariances with remote measured coordinates.  The
evolution of these covariances involves higher centered correlations.

\item Quantum empirical-measure methods reorganize this cascade, but they do
not close it at a fixed order.  Passing to the unnormalized equation, namely
the linear Zakai equation, removes the explicit normalization, but the exterior
observation records still carry information about the traced particles.

\item Entropy estimates typically require faithfulness or uniform absolute-continuity
conditions, which are not natural for the states considered here
and are not stable under the pure-state regimes relevant to the counting
approach.
\end{itemize}

There is a specific obstruction to the wave-function counting method.  If
\(
  p=|\psi\rangle\langle\psi|\in\cS(\HH)
\)
is a rank-one projection,  then
\begin{equation}
  \|\rho-p\|_1^2
  \le
  4\big(1-\Tr(\rho p)\big),
  \qquad \rho\in\cS(\HH).
  \label{eq:pure-state-trace-bound}
\end{equation}
This inequality allows one to control trace-norm errors through a scalar,
linear counting functional measuring the number of particles outside a
prescribed one-particle wave function.  It is precisely this pure-state
structure which underlies the wave-function approach.

Kolokoltsov extended the analysis of propagation of chaos to the stochastic setting via this counting functional in
\citep{kolokoltsov2021,kolokoltsov2022} under the assumptions
\[
  \operatorname{rank}(\rho_0^N)=1
  \qquad \mbox{and}\qquad 
  \eta=1,
\]
where  \(\eta=1\) ensures that $\operatorname{rank}(\rho_t^N)=1$ for all $t>0$ if $ \operatorname{rank}(\rho_0^N)=1$.   By contrast, when
\(0<\eta<1\), part of the environmental information is unobserved, and even a
pure initial state typically evolves into a mixed conditional state.  The
wave-function representation is then no longer available, and propagation of
chaos for general density matrices remained open, to the best of our knowledge.
Theorem~\ref{thm:main-approx} gives an affirmative answer for every efficiency
\(\eta\in(0,1]\) and arbitrary mixed one-particle initial data under the strong
tensorization assumption \eqref{eq:strong-tensorization}.

The proof uses three reductions which may also be useful for interacting
stochastic equations on density matrices.  First, at perfect efficiency
\(\eta=1\), mixed one-particle states are handled by purification: a mixed
state is realized as the partial trace of a pure state on an enlarged Hilbert
space,  a standard idea in quantum information theory, see e.g. Nielsen--Chuang
\cite{nielsenchuang2010} and Watrous \cite{watrous2018}.  In the present
many-particle setting, this requires lifting the Hamiltonian, interaction and
measurement operators and checking that the lifted Belavkin dynamics traces
down exactly to the original one.  Second, inefficient measurements are treated
by a fully observed dilation, followed by conditioning on the actually observed
channels and a relative-entropy argument which restores conditional
factorization.  Third, strong perturbations of the initial \(N\)-particle
state are controlled through the  Zakai equation, whose positivity and
trace preservation in expectation give an \(N\)-uniform stability estimate.
These ingredients avoid both the pure-state restriction and the closure
difficulties caused by exterior observation noises.

Finally, it is worth noting that an important exceptional regime occurs when
all measurement operators are skew-adjoint,  i.e.,
\[
  (L^{(k)})^*=-L^{(k)},
  \qquad
  k=1,\ldots,m.
\]
Hence the exterior Brownian motions disappear from the marginal equation.  A
stochastic BBGKY argument then becomes available.  

\medskip

For the sake of presentation, we shall use the following shorthand assumptions
to organize the proofs.

\begin{hypothesis}\label{hyp:perfect-efficiency-tensorization-purity}
The following three reductions will be used below.
\begin{enumerate}
\item Perfect tensorization \(\rho_0^N=\gamma_0^{\otimes N}\), perfect
efficiency \(\eta=1\), and quantum purity
\(\operatorname{rank}(\gamma_0)=1\).

\item Perfect tensorization \(\rho_0^N=\gamma_0^{\otimes N}\) and perfect
efficiency \(\eta=1\).

\item Perfect tensorization \(\rho_0^N=\gamma_0^{\otimes N}\).
\end{enumerate}
\end{hypothesis}








 




\section[Perfect efficiency and perfect tensorization]{Propagation of chaos under Hypothesis ~\ref{hyp:perfect-efficiency-tensorization-purity} (2): purification reduction}
\label{sec:purification}


Let Hypothesis~\ref{hyp:perfect-efficiency-tensorization-purity} (2) hold throughout Section \ref{sec:purification},  namely,  
\[
\rho_0^N=\gamma_0^{\otimes N} \qquad \mbox{and}\qquad  \eta=1.
\]
The goal of this section is to show the propagation of chaos using a purification reduction,  as summarized in Proposition \ref{prop:main-exact-eta-one}\footnote{The propagation of chaos result under Hypothesis~\ref{hyp:perfect-efficiency-tensorization-purity} (2) can be extended to infinite-dimensional Hilbert spaces $\HH$.} below. 


\begin{proposition}
\label{prop:main-exact-eta-one}
Let Hypothesis~\ref{hyp:perfect-efficiency-tensorization-purity} (2) hold with
$\gamma_0\in\cS(\HH)$.  
Then, for every fixed $n\ge1$ and $T>0$, there exists 
$C_{T,n}>0$, independent of $N$, such that
\begin{equation}
  \sup_{0\le t\le T}
  \E\left[
    \left\|
      \rho_t^{N:n}-\Gamma_t^n
    \right\|_1
  \right]
  \le
  C_{T,n}N^{-1/4}.
  \label{eq:eta-one-mixed-rate}
\end{equation}
\end{proposition}

The argument is a transfer principle.  Propagation of chaos is already known
when $\gamma_0$ is pure.  For a mixed $\gamma_0$, we represent it as the
partial trace of a pure state on an enlarged one-particle Hilbert space.  The
Hamiltonian, interaction and measurement operators are lifted by letting them
act trivially on the auxiliary factor.  The lifted dynamics is then a
perfectly observed pure-state dynamics, so the known theorem of Kolokoltsov \citep{kolokoltsov2021} applies.  The
physical mixed dynamics is recovered exactly by partial trace, and the desired
estimate follows from trace-norm contractivity.

We first record the convergence result of Kolokoltsov's pure-state
propagation theorem in the notation of the present paper,  see \citep{kolokoltsov2021} for related details.

\begin{theorem}[Kolokoltsov, 2021]
\label{thm:pure-input}
Let Hypothesis~\ref{hyp:perfect-efficiency-tensorization-purity} (1) hold.  
For $j\in[N]$, set
\begin{equation}
\rho_t^{N:(j)}
  :=
  \Tr_{[N]\setminus\{j\}}(\rho_t^N) \qquad \mbox{ and } \qquad  \alpha_{N,j}(t)
  :=
  1-\Tr\bigl(\rho_t^{N:(j)}\gamma_t^j\bigr).
  \label{eq:pure-counting-functional}
\end{equation}
Then, for every $T>0$, there exists a constant $C_T<\infty$, independent
of $N$ and $j$, such that
\begin{equation}
  \sup_{0\le t\le T}
  \E\bigl[\alpha_{N,j}(t)\bigr]
  \le
  \frac{C_T}{\sqrt N}.
  \label{eq:pure-counting-rate}
\end{equation}
Consequently, for every fixed $n\ge1$,
\begin{equation}
  \sup_{0\le t\le T}
  \E\left[
    \left\|
      \rho_t^{N:n}-\Gamma_t^n
    \right\|_1
  \right]
  \le
  2\sqrt{nC_T}\,N^{-1/4}.
  \label{eq:pure-trace-norm-rate}
\end{equation}
\end{theorem}

\begin{proof}
Estimate \eqref{eq:pure-counting-rate} is the finite-dimensional,
uncontrolled specialization of Kolokoltsov
\cite[Theorem~3.1]{kolokoltsov2021}.  The interaction \(V\) is symmetric by
assumption, and all coefficients are bounded in the present finite-dimensional
setting.  The same reference observes that the argument is unchanged for
\(m\ge1\) measurement operators
\(
  L^{(1)},\ldots,L^{(m)} .
\)

It remains to deduce \eqref{eq:pure-trace-norm-rate}.  Since $\eta=1$ and
$\gamma_0$ is pure, every $\gamma_t^j$ remains a rank-one projection.
For $j\le n$, denote by $P_{j,t}$ the lift of $\gamma_t^j$ to the
$j$-th factor of $\HH_N$.  The projections $P_{1,t},\ldots,P_{n,t}$
commute, and hence
\begin{equation}
  I_{\HH_N}
  -
  \prod_{j=1}^nP_{j,t}
  \le
  \sum_{j=1}^n
  \bigl(I_{\HH_N}-P_{j,t}\bigr),
  \label{eq:commuting-projection-bound}
\end{equation}
which yields
\begin{align}\label{eq:product-counting-bound}
  1-\Tr\bigl(\rho_t^{N:n}\Gamma_t^n\bigr)
  =
  \Tr\left(
    \rho_t^N
    \left(
      I_{\HH_N}-\prod_{j=1}^nP_{j,t}
    \right)
  \right)
 \le
  \sum_{j=1}^n
  \Tr\left(
    \rho_t^N
    \bigl(I_{\HH_N}-P_{j,t}\bigr)
  \right)=
  \sum_{j=1}^n\alpha_{N,j}(t).
\end{align}
As the product $\Gamma_t^n$ is also pure,  it follows that 
\[
  \left\|
    \rho_t^{N:n}-\Gamma_t^n
  \right\|_1
  \le
  2\sqrt{
    1-\Tr\bigl(\rho_t^{N:n}\Gamma_t^n\bigr)
  }
  \le
  2\left(
    \sum_{j=1}^n\alpha_{N,j}(t)
  \right)^{1/2}.
\]
Taking expectations, using Jensen's inequality and then
\eqref{eq:pure-counting-rate}, we obtain
\[
\begin{aligned}
  \E\left[
    \left\|
      \rho_t^{N:n}-\Gamma_t^n
    \right\|_1
  \right]
  \le
  2
  \left(
    \sum_{j=1}^n
    \E[\alpha_{N,j}(t)]
  \right)^{1/2}
  \le
  2\sqrt{nC_T}\,N^{-1/4}.
\end{aligned}
\]
Taking the supremum over $t\in[0,T]$ proves
\eqref{eq:pure-trace-norm-rate}.
\end{proof}

\begin{remark}
The rate $N^{-1/4}$ arises from the combination of the
$N^{-1/2}$ estimate for the counting functional and the square-root
comparison between that functional and the trace norm.  Optimality of this trace-norm rate remains open.  
\end{remark}

Next,  we adopt the purification approach and verify the lifted dynamics to remove the quantum purity assumption on $\gamma_0$.

\begin{lemma}
\label{lem:purification}
For every $\gamma_0\in\cS(\HH)$, there exist a finite-dimensional
auxiliary Hilbert space $\KK$, with
\(
  \dim(\KK)\le\dim(\HH),
\)
and a unit vector
\(
  \Psi_0\in\widetilde\HH:=\HH\otimes\KK
\)
such that, with
\(
  \widetilde\gamma_0
  :=
  |\Psi_0\rangle\langle\Psi_0|,
\)
one has
\begin{equation}
  \Tr_{\KK}(\widetilde\gamma_0)
  =
  \gamma_0.
  \label{eq:initial-purification}
\end{equation}
\end{lemma}

\begin{proof}
Let
\[
  \gamma_0
  =
  \sum_{r=1}^{r_0}
  \lambda_r|u_r\rangle\langle u_r|,
  \qquad \mbox{ with }
  r_0:=\operatorname{rank}(\gamma_0)
\]
be the spectral decomposition of $\gamma_0$, where $ (u_r)_{1\le r\le r_0}\subset \HH$ are orthogonal unit vectors and
\[
\lambda_r>0,
  \qquad
  \sum_{r=1}^{r_0}\lambda_r=1.
\]
Let $\KK$ be an $r_0$-dimensional Hilbert space with orthonormal basis
$(e_r)_{r=1}^{r_0}$.  Define
\[
  \Psi_0
  :=
  \sum_{r=1}^{r_0}
  \sqrt{\lambda_r}\,u_r\otimes e_r.
\]
Then $\|\Psi_0\|_{\widetilde\HH}=1$, and a straightforward computation yields
\[
  \widetilde\gamma_0
  =
  \sum_{r,s=1}^{r_0}
  \sqrt{\lambda_r\lambda_s}\,
  |u_r\rangle\langle u_s|
  \otimes
  |e_r\rangle\langle e_s|.
\]
Taking the partial trace over $\KK$ gives
\[
\begin{aligned}
  \Tr_{\KK}(\widetilde\gamma_0)
  &=
  \sum_{r,s=1}^{r_0}
  \sqrt{\lambda_r\lambda_s}\,
  |u_r\rangle\langle u_s|
  \langle e_s,e_r\rangle_{\KK}
  =
  \sum_{r=1}^{r_0}
  \lambda_r|u_r\rangle\langle u_r|
  =
  \gamma_0.
\end{aligned}
\]
This proves the lemma.
\end{proof}

Set
\(
  \widetilde\HH:=\HH\otimes\KK
\)
and define the lifted one-particle operators by
\begin{equation}
  \widetilde h
  :=
  h\otimes I_{\KK},
  \qquad
  \widetilde L^{(k)}
  :=
  L^{(k)}\otimes I_{\KK},
  \qquad
  k=1,\ldots,  m.
  \label{eq:lifted-one-particle-operators}
\end{equation}
Under the canonical identification
\(
  (\HH\otimes\KK)\otimes(\HH\otimes\KK)
  \simeq
  (\HH\otimes\HH)\otimes(\KK\otimes\KK),
\)
define
\begin{equation}
  \widetilde V
  :=
  V\otimes I_{\KK\otimes\KK}.
  \label{eq:lifted-interaction}
\end{equation}
Thus $\widetilde V$ acts as $V$ on the two physical factors and trivially
on the two auxiliary factors.  Define further 
\(
  \widetilde\HH_N
  :=
  \widetilde\HH^{\otimes N}
\)
and
\begin{equation}
  \widetilde H^N
  :=
  \sum_{j=1}^N\widetilde h_j
  +
  \frac1N
  \sum_{1\le j<\ell\le N}
  \widetilde V_{j\ell}.
  \label{eq:lifted-HN}
\end{equation}
Let
\[
  \mathbf T_N
  :=
  \Tr_{\KK^{\otimes N}}
  :
  \cL(\widetilde\HH_N)
  \longrightarrow
  \cL(\HH_N)
\]
denote the partial trace over all auxiliary factors.  The next lemma expresses the compatibility of the lifted coefficients with
the physical dynamics.

\begin{lemma}
\label{lem:trace-down}
For every $X\in\cL(\widetilde\HH_N)$,
$j\in[N]$ and $1\le k\le m$,
\begin{align}
  \mathbf T_N\left(
    -\i[\widetilde H^N,X]
  \right)
  &=
  -\i[H^N,\mathbf T_N(X)],
  \label{eq:trace-down-H}
  \\
  \mathbf T_N\left(
    \cD_{\widetilde L_j^{(k)}}(X)
  \right)
  &=
  \cD_{L_j^{(k)}}\bigl(\mathbf T_N(X)\bigr),
  \label{eq:trace-down-D}
  \\
  \mathbf T_N\left(
    \cC_{\widetilde L_j^{(k)}}(X)
  \right)
  &=
  \cC_{L_j^{(k)}}\bigl(\mathbf T_N(X)\bigr),
  \label{eq:trace-down-C}
  \\
  \mathbf T_N\left(
    \cE_{\widetilde L_j^{(k)}}(X)
  \right)
  &=
  \cE_{L_j^{(k)}}\bigl(\mathbf T_N(X)\bigr).
  \label{eq:trace-down-E}
\end{align}
Moreover, if
\(
  \widetilde\sigma\in\cS(\widetilde\HH)\) and \(
  \sigma:=\Tr_{\KK}(\widetilde\sigma),
\) 
then
\begin{equation}
  \widetilde V^{\widetilde\sigma}
  =
  V^\sigma\otimes I_{\KK}.
  \label{eq:trace-down-effective-potential}
\end{equation}
\end{lemma}

\begin{proof}
Under the canonical identification
\(
  \widetilde\HH_N
  \simeq
  \HH_N\otimes\KK^{\otimes N},
\)
the lifted Hamiltonian satisfies
\(
  \widetilde H^N
  =
  H^N\otimes I_{\KK^{\otimes N}}.
\)
For every $A\in\cL(\HH_N)$,
\begin{equation}
\begin{aligned}
  \mathbf T_N\left(
    (A\otimes I_{\KK^{\otimes N}})X
  \right)
  =
  A\mathbf T_N(X),\qquad 
  \mathbf T_N\left(
    X(A\otimes I_{\KK^{\otimes N}})
  \right)
  =
  \mathbf T_N(X)A.
\end{aligned}
  \label{eq:partial-trace-pull-through}
\end{equation}
Identity \eqref{eq:trace-down-H} follows immediately.

Since
\(
  \widetilde L_j^{(k)}
  =
  L_j^{(k)}\otimes I_{\KK^{\otimes N}},
\)
the same pull-through identities give
\[
\begin{aligned}
  \mathbf T_N\left(
    \widetilde L_j^{(k)}
    X
    (\widetilde L_j^{(k)})^*
  \right)
  &=
  L_j^{(k)}
  \mathbf T_N(X)
  (L_j^{(k)})^*,
  \\
  \mathbf T_N\left(
    (\widetilde L_j^{(k)})^*
    \widetilde L_j^{(k)}X
  \right)
  &=
  (L_j^{(k)})^*L_j^{(k)}
  \mathbf T_N(X),
  \\
  \mathbf T_N\left(
    X(\widetilde L_j^{(k)})^*
    \widetilde L_j^{(k)}
  \right)
  &=
  \mathbf T_N(X)
  (L_j^{(k)})^*L_j^{(k)}.
\end{aligned}
\]
This proves \eqref{eq:trace-down-D},  and
\eqref{eq:trace-down-C} follows in the same way.  Furthermore,
\[
\begin{aligned}
  \Tr\left(
    \bigl(
      \widetilde L_j^{(k)}
      +
      (\widetilde L_j^{(k)})^*
    \bigr)X
  \right)
  =
  \Tr\left(
    \bigl(
      L_j^{(k)}
      +
      (L_j^{(k)})^*
    \bigr)
    \mathbf T_N(X)
  \right).
\end{aligned}
\]
Combining this identity with \eqref{eq:trace-down-C} proves
\eqref{eq:trace-down-E}.

Finally, \eqref{eq:trace-down-effective-potential} follows from the
defining duality of the partial trace.  Indeed, for
$A\in\cL(\HH)$ and $B\in\cL(\KK)$,
\[
\begin{aligned}
  \Tr\left(
    (A\otimes B)
    \widetilde V^{\widetilde\sigma}
  \right)
  =
  \Tr(B)\,
  \Tr\left(
    (A\otimes\sigma)V
  \right)
  =
  \Tr(B)\,
  \Tr\left(
    AV^\sigma
  \right)
  =
  \Tr\left(
    (A\otimes B)
    (V^\sigma\otimes I_{\KK})
  \right).
\end{aligned}
\]
Operators of the form $A\otimes B$ span
$\cL(\widetilde\HH)$, which proves the identity.
\end{proof}

Let $\widetilde\rho^N$ solve the lifted $N$-particle Belavkin equation
\begin{equation}
\begin{aligned}
  \d\widetilde\rho_t^N
  &=
  -\i[
    \widetilde H^N,\widetilde\rho_t^N
  ]\,\d t
  +
  \sum_{k=1}^m\sum_{j=1}^N
  \cD_{\widetilde L_j^{(k)}}
  (\widetilde\rho_t^N)\,\d t
  +
  \sum_{k=1}^m\sum_{j=1}^N
  \cE_{\widetilde L_j^{(k)}}
  (\widetilde\rho_t^N)\,\d W_t^{k,j},
  \\
  \widetilde\rho_0^N
  &=
  \widetilde\gamma_0^{\otimes N}.
\end{aligned}
  \label{eq:lifted-N-eta-one}
\end{equation}
Let $\widetilde\gamma^j$ solve the lifted limiting equation
\begin{equation}
\begin{aligned}
  \d\widetilde\gamma_t^j
  &=
  \left(
    -\i[
      \widetilde h
      +
      \widetilde V^{\widetilde\xi_t},
      \widetilde\gamma_t^j
    ]
    +
    \sum_{k=1}^m
    \cD_{\widetilde L^{(k)}}
    (\widetilde\gamma_t^j)
  \right)\d t
  +
  \sum_{k=1}^m
  \cE_{\widetilde L^{(k)}}
  (\widetilde\gamma_t^j)\,\d W_t^{k,j},
  \\
  \widetilde\gamma_0^j
  &=
  \widetilde\gamma_0,
  \qquad
  \widetilde\xi_t
  :=
  \E[\widetilde\gamma_t^j].
\end{aligned}
  \label{eq:lifted-MF-eta-one}
\end{equation}
The initial states in both lifted equations \eqref{eq:lifted-N-eta-one} and \eqref{eq:lifted-MF-eta-one} are pure.  Since every output
channel is perfectly observed, the lifted equations preserve purity.  The next lemma shows that, after taking the partial trace, the lifted equations  reproduce the solutions to  \eqref{eq:N-Belavkin} and
\eqref{eq:MF-Belavkin}.   

\begin{lemma}
\label{lem:trace-down-dynamics}
The processes
\begin{equation}
  \rho_t^N
  :=
  \mathbf T_N(\widetilde\rho_t^N),
  \qquad
  \gamma_t^j
  :=
  \Tr_{\KK}(\widetilde\gamma_t^j)
  \label{eq:trace-down-processes}
\end{equation}
solve respectively \eqref{eq:N-Belavkin} and
\eqref{eq:MF-Belavkin} with $\eta=1$ and initial conditions
$\gamma_0^{\otimes N}$ and $\gamma_0$.
\end{lemma}

\begin{proof}
At time zero,
\[
\begin{aligned}
  \mathbf T_N
  \left(
    \widetilde\gamma_0^{\otimes N}
  \right)
  &=
  \left(
    \Tr_{\KK}(\widetilde\gamma_0)
  \right)^{\otimes N}
  =
  \gamma_0^{\otimes N}.
\end{aligned}
\]
Applying the bounded linear map $\mathbf T_N$ to the integral form of
\eqref{eq:lifted-N-eta-one} and using
Lemma~\ref{lem:trace-down} yields \eqref{eq:N-Belavkin} with
$\eta=1$.

For the limiting equation, set
\(
  \xi_t:=\E[\gamma_t^j].
\)
Linearity of expectation and partial trace gives
\[
  \Tr_{\KK}(\widetilde\xi_t)
  =\Tr_{\KK}\big( \E[\widetilde\gamma_t^j]\big)=
  \E\left[
    \Tr_{\KK}(\widetilde\gamma_t^j)
  \right]
  =  \E[\gamma_t^j]=
  \xi_t.
\]
Consequently, Lemma~\ref{lem:trace-down} gives
\[
  \widetilde V^{\widetilde\xi_t}
  =
  V^{\xi_t}\otimes I_{\KK}.
\]
Applying $\Tr_{\KK}$ to
\eqref{eq:lifted-MF-eta-one} therefore yields
\eqref{eq:MF-Belavkin} with $\eta=1$.  Pathwise uniqueness identifies
the trace-down processes with the solutions introduced in
Section~\ref{sec:framework}.
\end{proof}

We shall also use the following standard contractivity property.

\begin{lemma}
\label{lem:partial-trace-contraction}
For every $n\ge1$ and
$X\in\cL(\widetilde\HH^{\otimes n})$,
\begin{equation}
  \left\|
    \Tr_{\KK^{\otimes n}}(X)
  \right\|_1
  \le
  \|X\|_1.
  \label{eq:partial-trace-contraction}
\end{equation}
\end{lemma}

\begin{proof}
By trace duality,
\[
\begin{aligned}
  \left\|
    \Tr_{\KK^{\otimes n}}(X)
  \right\|_1
  =
  \sup_{A\in \cL(\HH^{\otimes n}):\, \|A\|\le1}
  \left|
    \Tr\left(
      A^*
      \Tr_{\KK^{\otimes n}}(X)
    \right)
  \right|=
  \sup_{A\in \cL(\HH^{\otimes n}):\,  \|A\|\le1}
  \left|
    \Tr\left(
      (A^*\otimes I_{\KK^{\otimes n}})X
    \right)
  \right|\le
  \|X\|_1,
\end{aligned}
\]
because
\(
  \|A\otimes I_{\KK^{\otimes n}}\|
  =
  \|A\|.
\)
\end{proof}

Now we collect all ingredients to prove Proposition~\ref{prop:main-exact-eta-one}.

\begin{proof}[Proof of Proposition~\ref{prop:main-exact-eta-one}]
For $1\le n\le N$, set
\[
  \widetilde\rho_t^{N:n}
  :=
  \Tr_{[N]\setminus[n]}(\widetilde\rho_t^N)
\qquad \mbox{ and }
\qquad 
  \widetilde\Gamma_t^n
  :=
  \widetilde\gamma_t^1
  \otimes\cdots\otimes
  \widetilde\gamma_t^n.
\]
Since $\widetilde\gamma_0$ is pure, Theorem~\ref{thm:pure-input},
applied on the enlarged one-particle space $\widetilde\HH$, gives a
constant $\widetilde C_{T,n}<\infty$, independent of $N$, such that
\begin{equation}
  \sup_{0\le t\le T}
  \E\left[
    \left\|
      \widetilde\rho_t^{N:n}
      -
      \widetilde\Gamma_t^n
    \right\|_1
  \right]
  \le
  \widetilde C_{T,n}N^{-1/4}.
  \label{eq:lifted-pure-chaos-rate}
\end{equation}
In view of Lemma~\ref{lem:trace-down-dynamics} and the associativity of partial
traces, it holds that
\(
  \rho_t^{N:n}
  =
  \Tr_{\KK^{\otimes n}}
  (\widetilde\rho_t^{N:n}).
\)
Moreover,
\[
\begin{aligned}
  \Gamma_t^n=
  \Tr_{\KK}(\widetilde\gamma_t^1)
  \otimes\cdots\otimes
  \Tr_{\KK}(\widetilde\gamma_t^n)=
  \Tr_{\KK^{\otimes n}}
  (\widetilde\Gamma_t^n),
\end{aligned}
\]
which gives
\[
  \rho_t^{N:n}-\Gamma_t^n
  =
  \Tr_{\KK^{\otimes n}}
  \left(
    \widetilde\rho_t^{N:n}
    -
    \widetilde\Gamma_t^n
  \right).
\]
Lemma~\ref{lem:partial-trace-contraction} therefore yields
\[
  \left\|
    \rho_t^{N:n}-\Gamma_t^n
  \right\|_1
  \le
  \left\|
    \widetilde\rho_t^{N:n}
    -
    \widetilde\Gamma_t^n
  \right\|_1,
\]
and proves \eqref{eq:eta-one-mixed-rate} by 
taking expectations and using
\eqref{eq:lifted-pure-chaos-rate}.
\end{proof}

\begin{remark}
The exact tensor-product assumption is used because
\[
  \gamma_0^{\otimes N}
  =
  \mathbf T_N
  \left(
    \widetilde\gamma_0^{\otimes N}
  \right),
\]
where $\widetilde\gamma_0^{\otimes N}$ is itself a pure tensor-product
state.  A general perturbation of $\gamma_0^{\otimes N}$ need not admit a
compatible pure tensor-product lifting.  Strong perturbations of the initial
condition are therefore treated separately in
Section~\ref{sec:initial-perturbation}.
\end{remark}


%%%%%%%%%%%%%%%%%%%%%%%%%%%%%%%%%%%%%%%%%%%%%%%%%%%%%%%%%%%%%%%%%%%%%%%%%%%%%%%
\section[Perfect tensorization: conditional and entropy approach]{Propagation of chaos under Hypothesis ~\ref{hyp:perfect-efficiency-tensorization-purity} (3):  conditional expectation and relative  entropy}
\label{sec:inefficient}


Let  $0<\eta<1$ and Hypothesis~\ref{hyp:perfect-efficiency-tensorization-purity} (3) hold in Section \ref{sec:inefficient},  namely, 
$$\rho^N_0=\gamma_0^{\otimes N}.$$
We aim to show the propagation of chaos using a conditional expectation argument,  as summarized in Proposition \ref{prop:main-exact-eta-less-one}. 


\begin{proposition}
\label{prop:main-exact-eta-less-one}
Let Hypothesis ~\ref{hyp:perfect-efficiency-tensorization-purity} (3) hold with
$\gamma_0\in\cS(\HH)$ and $\eta\in (0,1)$.    
Then, for every fixed $n\ge1$ and $T>0$, there exists 
$C_{\eta, T,n}>0$, independent of $N$, such that
\begin{equation}
  \sup_{0\le t\le T}
  \E\left[
    \left\|
      \rho_t^{N:n}-\Gamma_t^n
    \right\|_1
  \right]
  \le C_{\eta,T,n}N^{-1/16}.
  \label{eq:eta-main-rate}
\end{equation}
\end{proposition}

The proof uses a fully observed dilation.  For every  measurement
operator \(L^{(k)}\), introduce the two operators
\begin{equation}
  G^{(k,1)}:=\sqrt{\eta}\,L^{(k)},
  \qquad
  G^{(k,2)}:=\sqrt{1-\eta}\,L^{(k)},
  \qquad
  k=1,\ldots,  m.
  \label{eq:eta-dilated-operators}
\end{equation}
Since the coefficients are real and nonnegative, it follows that
\begin{align}
  &\cD_{G^{(k,1)}}+\cD_{G^{(k,2)}}
  =\cD_{L^{(k)}},
  \label{eq:eta-dilated-D}
  \\
  &\cE_{G^{(k,1)}}
  =\sqrt{\eta}\,\cE_{L^{(k)}},
  \qquad
  \cE_{G^{(k,2)}}
  =\sqrt{1-\eta}\,\cE_{L^{(k)}}.
  \label{eq:eta-dilated-E}
\end{align}
Thus the enlarged equation becomes a perfectly observed equation with \(2m\)
measurement operators.  Proposition~\ref{prop:main-exact-eta-one} applies to that equation.
The inefficient equation is then recovered by conditioning on the first
channel in each pair.  The nontrivial point is that conditioning does not
factorize at finite \(N\),  and this will be handled by a compensating Girsanov
transformation and a relative entropy estimate.




\subsection{Fully observed dilation and conditional projection}
\label{subsec:eta-dilation-projection}

For \(X\in\cL(\HH_N)\), set
\begin{equation}
  \mathfrak L_N(X)
  :=
  -\i[H^N,X]
  +
  \sum_{k=1}^m\sum_{j=1}^N
  \cD_{L_j^{(k)}}(X).
  \label{eq:eta-LN-def}
\end{equation}
Let \(\xi\) be the unique deterministic solution of
\eqref{eq:mean-state-equation}, and define the time-dependent one-particle
linear drift
\begin{equation}
  \mathfrak l_t(X)
  :=
  -\i[h+V^{\xi_t},X]
  +
  \sum_{k=1}^m
  \cD_{L^{(k)}}(X).
  \label{eq:eta-little-l-def}
\end{equation}
For later use, write
\begin{equation}
  A^{(k)}
  :=
  L^{(k)}+(L^{(k)})^*,
  \qquad
  b^{(k)}(x)
  :=
  \Tr\bigl(A^{(k)}x\bigr),
  \label{eq:eta-A-b-def}
\end{equation}
and, for \(R\in\cS(\HH_N)\),
\begin{equation}
  b_j^{N,k}(R)
  :=
  \Tr\bigl(A_j^{(k)}R\bigr)
  =
  \Tr\bigl(A^{(k)}R^{N:(j)}\bigr),
  \label{eq:eta-bNj-def}
\end{equation}
where
\(
  R^{N:(j)}
  :=
  \Tr_{[N]\setminus\{j\}}(R).
\)
Fix $T>0$.  Consider some filtered probability space
\(
  (\Omega,\cF,(\cF_t)_{0\le t\le T},\Q)
\)
which carries independent real-valued Brownian motions
\[
  Y^{k,j},\ U^{k,j},
  \qquad
  k=1,\ldots, m,
  \quad
   j=1,\ldots,  N.
\]
We write further
\(
  Y^j:=(Y^{1,j},\ldots,Y^{m,j})
\)
for the $m$-dimensional coordinate of particle $j$.

For $\vartheta\in\cL^+(\HH_N)$, let 
$\widehat R^{N,\vartheta}$ be the adapted process satisfying the Zakai equation below:
\begin{equation}
\begin{aligned}
 & \d\widehat R_t^{N,\vartheta}
  =
  \mathfrak L_N(\widehat R_t^{N,\vartheta})\,\d t
 +
  \sqrt{\eta}
  \sum_{k=1}^m\sum_{j=1}^N
  \cC_{L_j^{(k)}}(\widehat R_t^{N,\vartheta})\,\d Y_t^{k,j}
  +
  \sqrt{1-\eta}
  \sum_{k=1}^m\sum_{j=1}^N
  \cC_{L_j^{(k)}}(\widehat R_t^{N,\vartheta})\,\d U_t^{k,j},
  \\
  &\widehat R_0^{N,\vartheta}
  =\vartheta.
\end{aligned}
  \label{eq:eta-full-linear-N}
\end{equation}


\begin{lemma}
\label{lem:full-linear-reference}
For every $t\in[0,T]$,  define the random map
\[
  \Psi_t^{N,\eta}:\vartheta
  \longmapsto
  \widehat R_t^{N,\vartheta}.
\]
Then the map is linear and positive (almost surely).  Moreover, for every
$\vartheta\in\cL^+(\HH_N)$,
\begin{equation}
  \E_{\Q}\left[
    \Tr(\widehat R_t^{N,\vartheta})
  \right]
  =
  \Tr(\vartheta).
  \label{eq:eta-full-mean-trace}
\end{equation}
\end{lemma}

\begin{proof}
Let $\mathscr U^N$ be the unique solution to the following SDE taking values in $\cL(\HH_N)$:
\[
\begin{aligned}
  \d\mathscr U_t^N
  &=
  \left(
    -\i H^N
    -\frac12
    \sum_{k=1}^m\sum_{j=1}^N
    (L_j^{(k)})^*L_j^{(k)}
  \right)\mathscr U_t^N\,\d t+
  \sum_{k=1}^m\sum_{j=1}^N
  \left(
    \sqrt{\eta}\,L_j^{(k)}\,\d Y_t^{k,j}
    +
    \sqrt{1-\eta}\,L_j^{(k)}\,\d U_t^{k,j}
  \right)\mathscr U_t^N,
  \\
  \mathscr U_0^N&=I_{\HH_N}.
\end{aligned}
\]
As  \eqref{eq:eta-full-linear-N} admits a unique pathwise solution,  It\^o's formula gives,  
\begin{equation}
  \widehat R_t^{N,\vartheta}
  =
  \mathscr U_t^N\vartheta(\mathscr U_t^N)^*.
  \label{eq:eta-full-propagator}
\end{equation}
Hence the map is linear and  positive.  Taking the trace and
expectation in \eqref{eq:eta-full-linear-N} gives
\eqref{eq:eta-full-mean-trace}, since the Hamiltonian and Lindblad drifts
have zero trace and the stochastic integrals have zero expectation.
\end{proof}

We now specialize to
\(
  \vartheta=\gamma_0^{\otimes N}
\)
and omit the superscript $\vartheta$.  Set
\begin{equation}
  M_t^N:=\Tr(\widehat R_t^N)
  \qquad\mbox{ and } \qquad 
  R_t^N:=\frac{\widehat R_t^N}{M_t^N}.
  \label{eq:eta-MR-def}
\end{equation}
The matrix stochastic exponential $\mathscr U_t^N$ is invertible almost surely.  Taking the trace in
\eqref{eq:eta-full-linear-N} yields
\[
\begin{aligned}
  \d M_t^N
  &=
  M_t^N
  \sum_{k=1}^m\sum_{j=1}^N
  \bigg(
    \sqrt{\eta}\,b_j^{N,k}(R_t^N)\,\d Y_t^{k,j}+
    \sqrt{1-\eta}\,b_j^{N,k}(R_t^N)\,\d U_t^{k,j}
  \bigg).
\end{aligned}
\]
The functions $b_j^{N,k}$ are uniformly bounded on the state space, so
Novikov's condition holds.  Thus $M^N$ is a strictly positive martingale
with mean one, and  \eqref{eq:eta-full-propagator} implies $\widehat R_t^N\ge0$ and $R_t^N\in\cS(\HH_N)$.   In addition,
\begin{equation}
  \frac{\d\Q^{N,\mathrm{full}}}{\d\Q}
  \bigg|_{\cF_T}
  :=M_T^N
  \label{eq:eta-full-measure}
\end{equation}
defines a probability measure equivalent to $\Q$.  Under
$\Q^{N,\mathrm{full}}$, the processes
\begin{align}
  B_t^{k,j,1}
  &:={}
  Y_t^{k,j}
  -\sqrt{\eta}
  \int_0^t b_j^{N,k}(R_s^N)\,\d s,
  \label{eq:eta-B1-full}
  \\
  B_t^{k,j,2}
  &:={}
  U_t^{k,j}
  -\sqrt{1-\eta}
  \int_0^t b_j^{N,k}(R_s^N)\,\d s
  \label{eq:eta-B2-full}
\end{align}
form a family of  independent Brownian motions.  Normalizing
\eqref{eq:eta-full-linear-N} together with Itô's formula yields
\begin{equation}
\begin{aligned}
  \d R_t^N
  &=
  \mathfrak L_N(R_t^N)\,\d t
 +
  \sqrt{\eta}
  \sum_{k=1}^m\sum_{j=1}^N
  \cE_{L_j^{(k)}}(R_t^N)\,\d B_t^{k,j,1}+
  \sqrt{1-\eta}
  \sum_{k=1}^m\sum_{j=1}^N
  \cE_{L_j^{(k)}}(R_t^N)\,\d B_t^{k,j,2}.
\end{aligned}
  \label{eq:eta-full-nonlinear-N}
\end{equation}
On the same probability space,  let \(X^1,\ldots,X^N\) solve
\begin{equation}
\begin{aligned}
  \d X_t^j
  &=
  \mathfrak l_t(X_t^j)\,\d t
  +
  \sqrt{\eta}
  \sum_{k=1}^m
  \cE_{L^{(k)}}(X_t^j)\,\d B_t^{k,j,1}
  +
  \sqrt{1-\eta}
  \sum_{k=1}^m
  \cE_{L^{(k)}}(X_t^j)\,\d B_t^{k,j,2},
  \\
  X_0^j&=\gamma_0,
\end{aligned}
  \label{eq:eta-full-MF}
\end{equation}
where we note that  \(X^1,\ldots,X^N\) are well defined under both probability measures $\Q$,  
$\Q^{N,\mathrm{full}}$.  
The processes \(X^1,\ldots,X^N\) are independent and identically
distributed under \(\Q^{N,\mathrm{full}}\).  Taking expectations in
\eqref{eq:eta-full-MF},  one finds that  $t\longmapsto  \E_{\Q^{N,\mathrm{full}}}[X_t^j]$ solves \eqref{eq:mean-state-equation},  which shows by means of the uniqueness for
\eqref{eq:mean-state-equation} 
\begin{equation}
  \E_{\Q^{N,\mathrm{full}}}[X_t^j]
  =
  \xi_t.
  \label{eq:eta-full-MF-mean}
\end{equation}

\begin{lemma}
\label{lem:eta-full-exchangeability}
Under $\Q^{N,\mathrm{full}}$, the family
\(
  \bigl(R^N,(X^j,Y^j)_{j=1}^N\bigr)
\)
is exchangeable.  In particular, the path-space law
of the pairs $(X^j,Y^j)$ is invariant under every permutation of the
particle labels.
\end{lemma}

\begin{proof}
Fix $\pi\in\mathfrak S_N$.  Conjugating the state by $U_\pi$ and replacing
$(B^{k,j,a})_{k,j,a}$ with
$(B^{k,\pi^{-1}(j),a})_{k,j,a}$ leaves the coupled equations
\eqref{eq:eta-full-nonlinear-N} and \eqref{eq:eta-full-MF} unchanged,
because $H^N$ is permutation invariant, all particles have the same
measurement operators, and the initial state is
$\gamma_0^{\otimes N}$.  The permuted Brownian family has the same law as
the original one.  Uniqueness in law for the coupled system therefore gives
the asserted equality in law.  The observation coordinates are recovered
from
\[
  \d Y_t^{k,j}
  =
  \sqrt{\eta}\,b_j^{N,k}(R_t^N)\,\d t
  +\d B_t^{k,j,1},
\]
so they are permuted simultaneously with the corresponding particles.
\end{proof}

\begin{proposition}
\label{prop:eta-full-chaos}
For every fixed \(n\ge1\) and \(T>0\), there exists
\(C_{T,n}^{\mathrm{full}}<\infty\), independent of \(N\), such that
\begin{equation}
\begin{aligned}
  \varepsilon_{N,n}^{\mathrm{full}}(T)
  &:={}
  \sup_{0\le t\le T}
  \E_{\Q^{N,\mathrm{full}}}
  \left[
    \left\|
      R_t^{N:n}
      -
      X_t^1\otimes\cdots\otimes X_t^n
    \right\|_1
  \right] \le
  C_{T,n}^{\mathrm{full}}N^{-1/4}.
\end{aligned}
  \label{eq:eta-full-chaos}
\end{equation}
\end{proposition}

\begin{proof}
In view of \eqref{eq:eta-dilated-D}--\eqref{eq:eta-dilated-E}, equation
\eqref{eq:eta-full-nonlinear-N} is exactly the perfect-efficiency Belavkin equation
with the \(2m\) measurement operators
\[
  G^{(k,1)},\ G^{(k,2)},
  \qquad
  k=1,\ldots, m, 
\]
and equation \eqref{eq:eta-full-MF} is the corresponding mean-field equation.
Proposition~\ref{prop:main-exact-eta-one}, applied to this finite family of
operators, gives \eqref{eq:eta-full-chaos}.
\end{proof}

We now retain only the observed outputs \(Y^{k,j}\).  Define the filtration 
\begin{equation}
  \cY_t^N
  :=
  \sigma\bigl(
    Y_s^{k,j}:
    0\le s\le t,\ 1\le k\le m,\ 1\le j\le N
  \bigr),
  \label{eq:eta-YN-def}
\end{equation}
with the usual completion, and define the conditional expectation
\begin{equation}
  \widehat\rho_t^{N,\eta}
  :=
  \E_{\Q}
  \left[
    \widehat R_t^N\mid\cY_t^N
  \right].
  \label{eq:eta-hat-rho-conditional}
\end{equation}

\begin{proposition}
\label{prop:eta-conditional-reduction}
The process \(\widehat\rho^{N,\eta}\) defined by  \eqref{eq:eta-hat-rho-conditional} must satisfy
\begin{equation}
\begin{aligned}
  \d\widehat\rho_t^{N,\eta}
  &=
  \mathfrak L_N(\widehat\rho_t^{N,\eta})\,\d t+
  \sqrt{\eta}
  \sum_{k=1}^m\sum_{j=1}^N
  \cC_{L_j^{(k)}}(\widehat\rho_t^{N,\eta})\,\d Y_t^{k,j},
  \\
  \widehat\rho_0^{N,\eta}
  &=
  \gamma_0^{\otimes N}.
\end{aligned}
  \label{eq:eta-one-channel-linear}
\end{equation}
Set
\begin{equation}
  Z_t^{N,\eta}
  :=
  \Tr(\widehat\rho_t^{N,\eta})>0,
  \qquad
  \rho_t^{N,\eta}
  :=
  \frac{\widehat\rho_t^{N,\eta}}{Z_t^{N,\eta}}.
  \label{eq:eta-Z-rho-def}
\end{equation}
Then
\begin{equation}
  Z_t^{N,\eta}
  =
  \E_{\Q}
  \left[
    M_t^N\mid\cY_t^N
  \right].
  \label{eq:eta-Z-as-conditional-M}
\end{equation}
The measure $\Q^{N,\eta}$ given by
\begin{equation}
  \frac{\d\Q^{N,\eta}}{\d\Q}
  \bigg|_{\cY_T^N}
  :=
  Z_T^{N,\eta}
  \label{eq:eta-Qeta-def}
\end{equation}
is the restriction of \(\Q^{N,\mathrm{full}}\) to \(\cY_T^N\), and it holds that
\begin{equation}
  \rho_t^{N,\eta}
  =
  \E_{\Q^{N,\mathrm{full}}}
  \left[
    R_t^N\mid\cY_t^N
  \right].
  \label{eq:eta-normalized-conditional}
\end{equation}
Moreover,  the measures \(\Q^{N,\eta}\) and
\(\Q\) are equivalent on \(\cY_T^N\), and the processes
\begin{equation}
  W_t^{N,k,j}
  :=
  Y_t^{k,j}
  -
  \sqrt{\eta}
  \int_0^t
  b_j^{N,k}(\rho_s^{N,\eta})\,\d s
  \label{eq:eta-WN-def}
\end{equation}
form a family of  independent Brownian motions under \(\Q^{N,\eta}\).  In particular,  
the normalized process solves
\begin{equation}
\begin{aligned}
  \d\rho_t^{N,\eta}
  &=
  \mathfrak L_N(\rho_t^{N,\eta})\,\d t+
  \sqrt{\eta}
  \sum_{k=1}^m\sum_{j=1}^N
  \cE_{L_j^{(k)}}(\rho_t^{N,\eta})\,\d W_t^{N,k,j}.
\end{aligned}
  \label{eq:eta-N-target}
\end{equation}
\end{proposition}

\begin{proof}
We first justify the projection of the stochastic integrals.  If \(F\) is a
square-integrable \((\cF_t)\)-predictable process, then, coordinatewise in
\(\cL(\HH_N)\),
\begin{equation}
  \E_{\Q}
  \left[
    \int_0^t F_s\,\d U_s^{k,j}
    \mathrel{\bigg|}
    \cY_t^N
  \right]
  =0.
  \label{eq:eta-U-projection}
\end{equation}
Indeed, for an elementary predictable process and a bounded
\(\cY_t^N\)-measurable random variable \(H\), each expectation is a sum of
terms of the form
\[
  \E_{\Q}
  \left[
    H F_{t_\ell}
    \bigl(
      U_{t_{\ell+1}}^{k,j}-U_{t_\ell}^{k,j}
    \bigr)
  \right],
\]
which vanish because the displayed increment is centered and independent of
the entire \(Y\)-path and of the joint past at time \(t_\ell\).  The general
case follows by the It\^o isometry.  Notice that the integrand may depend on
both \(Y\) and \(U\).

Similarly,
\begin{equation}
  \E_{\Q}
  \left[
    \int_0^t F_s\,\d Y_s^{k,j}
    \mathrel{\bigg|}
    \cY_t^N
  \right]
  =
  \int_0^t
  \E_{\Q}
  \left[
    F_s\mid\cY_s^N
  \right]\d Y_s^{k,j}.
  \label{eq:eta-Y-projection}
\end{equation}
For \(s\le t\), the future increments of \(Y\) after time \(s\) are
independent of \(\cF_s\), so
\[
  \E_{\Q}
  \left[
    \widehat R_s^N\mid\cY_t^N
  \right]
  =
  \E_{\Q}
  \left[
    \widehat R_s^N\mid\cY_s^N
  \right].
\]
Applying these identities to the integral form of
\eqref{eq:eta-full-linear-N}, and using the linearity of all coefficients,
gives \eqref{eq:eta-one-channel-linear}.

Equation \eqref{eq:eta-Z-as-conditional-M} follows by taking the trace.  For
every \(A\in\cY_T^N\),
\[
\begin{aligned}
  \Q^{N,\mathrm{full}}(A)
  =
  \E_{\Q}[\1_A M_T^N]
  =
  \E_{\Q}
  \left[
    \1_A
    \E_{\Q}[M_T^N\mid\cY_T^N]
  \right]
  =
  \E_{\Q}[\1_A Z_T^{N,\eta}],
\end{aligned}
\]
which proves that \(\Q^{N,\eta}\) is the stated restriction.

Bayes' formula gives
\[
\begin{aligned}
  \E_{\Q^{N,\mathrm{full}}}
  \left[
    R_t^N\mid\cY_t^N
  \right]
  &=
  \frac{
    \E_{\Q}
    \left[
      M_T^N R_t^N\mid\cY_t^N
    \right]
  }{
    \E_{\Q}
    \left[
      M_T^N\mid\cY_t^N
    \right]
  }
  =
  \frac{
    \E_{\Q}
    \left[
      M_t^N R_t^N\mid\cY_t^N
    \right]
  }{Z_t^{N,\eta}}
  =
  \rho_t^{N,\eta},
\end{aligned}
\]
because \(R_t^N\) is \(\cF_t\)-measurable and \(M^N\) is a
\(\Q\)-martingale.  Since \(M_t^N>0\), its conditional expectation
\(Z_t^{N,\eta}\) is strictly positive.  Finally, taking the trace in
\eqref{eq:eta-one-channel-linear} shows that
\[
  \d Z_t^{N,\eta}
  =
  \sqrt{\eta}\,Z_t^{N,\eta}
  \sum_{k=1}^m\sum_{j=1}^N
  b_j^{N,k}(\rho_t^{N,\eta})\,\d Y_t^{k,j}.
\]
The integrands are bounded.  Girsanov's theorem therefore gives
\eqref{eq:eta-WN-def}, and normalization of
\eqref{eq:eta-one-channel-linear} with Itô's formula gives \eqref{eq:eta-N-target}.
\end{proof}

For \(1\le n\le N\), set
\begin{equation}
  \mathbf X_t^n
  :=
  X_t^1\otimes\cdots\otimes X_t^n,
  \qquad
  \Theta_t^{N,n}
  :=
  \E_{\Q^{N,\mathrm{full}}}
  \left[
    \mathbf X_t^n\mid\cY_t^N
  \right],
  \label{eq:eta-product-Theta-def}
\end{equation}
where we recall that \(X^1,\ldots,X^N\) are defined by 
  \eqref{eq:eta-full-MF}.  Since partial trace commutes with conditional expectation,  one has
\[
  \rho_t^{N:n,\eta}
  =
  \E_{\Q^{N,\mathrm{full}}}
  \left[
    R_t^{N:n}\mid\cY_t^N
  \right].
\]
Conditional Jensen's inequality and the equality of
\(\Q^{N,\eta}\) and \(\Q^{N,\mathrm{full}}\) on \(\cY_t^N\) give
\begin{equation}
\begin{aligned}
  &\sup_{0\le t\le T}
  \E_{\Q^{N,\eta}}
  \left[
    \left\|
      \rho_t^{N:n,\eta}-\Theta_t^{N,n}
    \right\|_1
  \right]
  \le
  C_{T,n}^{\mathrm{full}}N^{-1/4},
\end{aligned}
  \label{eq:eta-first-term-bound}
\end{equation}
for some constant $C_{T,n}^{\mathrm{full}}>0$ appearing in Proposition \ref{prop:eta-full-chaos}.


\subsection{Relative entropy and conditional factorization}
\label{subsec:eta-conditional-factorization}


For \(0\le t\le T\),  define the Polish spaces as follows
\[
  \mathsf X_t:=C([0,t],\cS(\HH)),
  \qquad
  \mathsf Y_t:=C([0,t],\mathbb R^m).
\]
Let \(\mathbf q_T\) be the law on
\(\mathsf X_T\times\mathsf Y_T\) of a pair \((X,Y)\) satisfying
\begin{equation}
\begin{aligned}
  \d X_t
  &=
  \mathfrak l_t(X_t)\,\d t
  +
  \sqrt{\eta}
  \sum_{k=1}^m
  \cE_{L^{(k)}}(X_t)\,\d\beta_t^{k,1}
  +
  \sqrt{1-\eta}
  \sum_{k=1}^m
  \cE_{L^{(k)}}(X_t)\,\d\beta_t^{k,2},
  \qquad
  X_0=\gamma_0,
  \\
  \d Y_t^k
  &=
  \sqrt{\eta}\,b^{(k)}(X_t)\,\d t
  +
  \d\beta_t^{k,1},
  \qquad
  Y_0^k=0,
\end{aligned}
  \label{eq:eta-ideal-pair}
\end{equation}
where \(\{\beta^{k,a}:1\le k\le m,\ a=1,2\}\) are independent
Brownian motions.  Let \(\mathbf P_{N,T}\) denote the joint law under
\(\Q^{N,\mathrm{full}}\) of
\(
  (X^1,Y^1),\ldots,(X^N,Y^N)
\) with $Y^j=(Y^{1,j},\ldots,Y^{m,j})$.  Define
\begin{equation}
  \Delta_t^{k,j}
  :=
  b_j^{N,k}(R_t^N)-b^{(k)}(X_t^j)
  \label{eq:eta-Delta-def}
\end{equation}
and
\begin{equation}
  u_t^{k,j,1}
  :=
  \sqrt{\eta}\,\Delta_t^{k,j},
  \qquad
  u_t^{k,j,2}
  :=
  -\frac{\eta}{\sqrt{1-\eta}}\,\Delta_t^{k,j}.
  \label{eq:eta-controls-def}
\end{equation}
As the controls are bounded and predictable,  define $\widetilde\Q^N$ by 
\begin{equation}
\begin{aligned}
  \frac{\d\widetilde\Q^N}{\d\Q^{N,\mathrm{full}}}
  \bigg|_{\cF_T}
  &:={}
  \exp\Bigg(
    -
    \sum_{k=1}^m\sum_{j=1}^N\sum_{a=1}^2
    \int_0^T u_s^{k,j,a}\,\d B_s^{k,j,a}    -
    \frac12
    \sum_{k=1}^m\sum_{j=1}^N\sum_{a=1}^2
    \int_0^T |u_s^{k,j,a}|^2\,\d s
  \Bigg).
\end{aligned}
  \label{eq:eta-compensating-density}
\end{equation}
Again,  Novikov's condition holds,  then under \(\widetilde\Q^N\), the processes
\begin{equation}
  \widetilde B_t^{k,j,a}
  :=
  B_t^{k,j,a}
  +
  \int_0^t u_s^{k,j,a}\,\d s
  \label{eq:eta-Btilde-def}
\end{equation}
form a family of independent Brownian motions,  and a straightforward verification gives
\[
  \d Y_t^{k,j}
  =
  \sqrt{\eta}\,b^{(k)}(X_t^j)\,\d t
  +
  \d\widetilde B_t^{k,j,1}.
\]
If we replace \(B^{k,j,a}\)  by \(\widetilde B^{k,j,a}\) in the SDE
for \(X^j\), the additional drift corresponding to the \(k\)-th operator becomes
\[
  -\cE_{L^{(k)}}(X_t^j)
  \left(
    \sqrt{\eta}\,u_t^{k,j,1}
    +
    \sqrt{1-\eta}\,u_t^{k,j,2}
  \right)
  =0.
\]
Hence, by pathwise uniqueness,
\begin{equation}
  \Law_{\widetilde\Q^N}
  \bigl(
    (X^1,Y^1),\ldots,(X^N,Y^N)
  \bigr)
  =
  \mathbf q_T^{\otimes N}.
  \label{eq:eta-product-law-under-tilde}
\end{equation}
For probability measures \(\mu,\nu\),  we define the relative entropy by
\[
  \Ent(\mu\mid\nu)
  :=
  \int
  \log\left(\frac{\d\mu}{\d\nu}\right)\d\mu
\]
when \(\mu\ll\nu\), and set the entropy equal to \(+\infty\) otherwise.
Set
\begin{equation}
  \Lambda_L
  :=
  \sum_{k=1}^m
  \|A^{(k)}\|^2.
  \label{eq:eta-LambdaL-def}
\end{equation}
Then we have the following control.
\begin{proposition}
\label{prop:eta-entropy}
For every \(T>0\),  it holds that
\begin{equation}
\begin{aligned}
  \frac1N
  \Ent\left(
    \mathbf P_{N,T}
    \mid
    \mathbf q_T^{\otimes N}
  \right)
  &\le
  \frac{\eta \Lambda_L}{1-\eta}
  \int_0^T
  \E_{\Q^{N,\mathrm{full}}}
  \left[
    \left\|
      R_t^{N:(1)}-X_t^1
    \right\|_1
  \right]\d t\le
  C_{\eta,T}N^{-1/4},
\end{aligned}
  \label{eq:eta-entropy-bound}
\end{equation}
for some $C_{\eta, T}>0$.
\end{proposition}

\begin{proof}
By data processing and \eqref{eq:eta-product-law-under-tilde},  it holds that 
\[
  \Ent\left(
    \mathbf P_{N,T}
    \mid
    \mathbf q_T^{\otimes N}
  \right)
  \le
  \Ent\left(
    \Q^{N,\mathrm{full}}
    \mid
    \widetilde\Q^N
  \right).
\]
Girsanov's entropy identity gives
\[
\begin{aligned}
  \Ent\left(
    \Q^{N,\mathrm{full}}
    \mid
    \widetilde\Q^N
  \right)
  =
  \frac12
  \sum_{k=1}^m\sum_{j=1}^N
  \E_{\Q^{N,\mathrm{full}}}
  \int_0^T
  \left(
    |u_t^{k,j,1}|^2+|u_t^{k,j,2}|^2
  \right)\d t
  =
  \frac{\eta}{2(1-\eta)}
  \sum_{k=1}^m\sum_{j=1}^N
  \E_{\Q^{N,\mathrm{full}}}
  \int_0^T
  |\Delta_t^{k,j}|^2\,\d t.
\end{aligned}
\]
By Lemma~\ref{lem:eta-full-exchangeability}, division by \(N\) reduces the sum over \(j\) to \(j=1\).  Moreover,
\[
  |\Delta_t^{k,1}|^2
  \le
  \|A^{(k)}\|^2
  \left\|
    R_t^{N:(1)}-X_t^1
  \right\|_1^2  \le
  2\|A^{(k)}\|^2
  \left\|
    R_t^{N:(1)}-X_t^1
  \right\|_1,  
\]
since the trace distance between density matrices is at most two.  
Summing over \(k\) proves the first inequality in
\eqref{eq:eta-entropy-bound}.  The second follows from
Proposition~\ref{prop:eta-full-chaos} with \(n=1\).
\end{proof}

We next recall two standard entropy facts.  If \(\mu\) and \(\nu\) are
probability measures on a product \(E\times F\) of Polish spaces, then
\begin{equation}
  \Ent(\mu\mid\nu)
  =
  \Ent(\mu_F\mid\nu_F)
  +
  \int_F
  \Ent\left(
    \mu^{E\mid F=y}
    \mid
    \nu^{E\mid F=y}
  \right)\mu_F(\d y).
  \label{eq:eta-entropy-chain-rule}
\end{equation}
Furthermore, if \(\mu\) is a probability measure on \(E^N\),
\(\nu_1,\ldots,\nu_N\) are probability measures on \(E\), and
\(I_1,\ldots,I_r\subset[N]\) are pairwise disjoint, then
\begin{equation}
  \sum_{\ell=1}^r
  \Ent\left(
    \mu_{I_\ell}
    \mathrel{\bigg|}
    \bigotimes_{j\in I_\ell}\nu_j
  \right) \le \Ent\left(
    \mu
    \mathrel{\bigg|}
    \bigotimes_{j=1}^N\nu_j
  \right).
  \label{eq:eta-block-superadditivity}
\end{equation}
The latter follows by applying the entropy chain rule successively to the
blocks and discarding the nonnegative mutual-information terms.

For \(t\le T\), let \(\mathbf P_{N,t}\) and \(\mathbf q_t\) be the
restrictions of \(\mathbf P_{N,T}\) and \(\mathbf q_T\) to \([0,t]\).
Let \(\mathbf P_{N,t}^{X\mid Y}\) be a regular conditional law of
\((X^1,\ldots,X^N)\) given \((Y^1,\ldots,Y^N)\), and let
\(\mathbf q_t^{X\mid Y=y}\) be a regular conditional law under
\(\mathbf q_t\).  The chain rule and monotonicity of entropy under
restriction lead to
\begin{equation}
\begin{aligned}
  &\E_{\mathbf P_{N,t}^Y}
  \Ent\left(
    \mathbf P_{N,t}^{X\mid Y}
    \mathrel{\bigg|}
    \bigotimes_{j=1}^N
    \mathbf q_t^{X\mid Y=Y^j}
  \right)\le
  \Ent\left(
    \mathbf P_{N,T}
    \mid
    \mathbf q_T^{\otimes N}
  \right).
\end{aligned}
  \label{eq:eta-total-conditional-entropy}
\end{equation}
Fix \(n\ge1\), set \(r_N:=\lfloor N/n\rfloor\), and partition the first \(nr_N\) coordinates into \(r_N\) consecutive blocks of size \(n\).  Applying \eqref{eq:eta-block-superadditivity} to the conditional law \(\mathbf P_{N,t}^{X\mid Y=y^{1:N}}\), with reference measure \(\bigotimes_{j=1}^N\mathbf q_t^{X\mid Y=y^j}\), and then integrating with respect to the \(Y\)-marginal \(\mathbf P_{N,t}^Y\),  i.e.,  \(\mathbf P_{N,t}^Y\) denotes the law of \((Y^1,\ldots,Y^N)\) under \(\mathbf P_{N,t}\),  gives a sum of \(r_N\) block entropies.  The joint law \(\mathbf P_{N,t}\) is invariant under simultaneous permutations of the pairs \((X^j,Y^j)\).  Hence,  the expectations of these block entropies are equal.  Therefore,
\begin{equation}
\begin{aligned}
  &\E_{\mathbf P_{N,t}^Y}
  \Ent\left(
    \mathbf P_{N,t}^{X^{1:n}\mid Y^{1:N}}
    \mathrel{\bigg|}
    \bigotimes_{j=1}^n
    \mathbf q_t^{X\mid Y=Y^j}
  \right) \le
  \frac1{r_N}
  \Ent\left(
    \mathbf P_{N,T}
    \mid
    \mathbf q_T^{\otimes N}
  \right).
\end{aligned}
  \label{eq:eta-block-conditional-entropy}
\end{equation}
Define the conditional barycenter
\begin{equation}
  \overline\gamma_t(y)
  :=
  \int_{\mathsf X_t}
  x_t\,
  \mathbf q_t^{X\mid Y=y}(\d x)
  \label{eq:eta-bar-gamma-barycenter}
\end{equation}
and
\begin{equation}
  \overline\Gamma_t^{N,n}
  :=
  \overline\gamma_t(Y^1)
  \otimes\cdots\otimes
  \overline\gamma_t(Y^n).
  \label{eq:eta-bar-Gamma-def}
\end{equation}
The conditional kernels may be chosen arbitrarily on a
\(\mathbf q_t^Y\)-null set.  This causes no ambiguity here because
\(\mathbf P_{N,t}^Y\ll(\mathbf q_t^Y)^{\otimes N}\), as follows from the
finite entropy in Proposition~\ref{prop:eta-entropy}.  A canonical version
will be selected in Lemma~\ref{lem:eta-barycenter-identification} below.

\begin{proposition}
\label{prop:eta-conditional-factorization}
For every fixed \(n\ge1\) and \(T>0\),  there exists $C_{\eta,T,n}>0$ so that
\begin{equation}
\begin{aligned}
  \sup_{0\le t\le T}
  \E_{\Q^{N,\eta}}
  \left[
    \left\|
      \Theta_t^{N,n}-\overline\Gamma_t^{N,n}
    \right\|_1
  \right]\le
  \left[
    \frac{2\eta N \Lambda_L}{(1-\eta)r_N}
    \int_0^T
    \E_{\Q^{N,\mathrm{full}}}
    \left[
      \left\|
        R_s^{N:(1)}-X_s^1
      \right\|_1
    \right]\d s
  \right]^{1/2}
\le
  C_{\eta,T,n}N^{-1/8}.
\end{aligned}
  \label{eq:eta-conditional-factorization-bound}
\end{equation}
\end{proposition}

\begin{proof}
Conditionally on \(Y^{1:N}\), Pinsker's inequality and trace duality give
\[
\begin{aligned}
  \left\|
    \Theta_t^{N,n}-\overline\Gamma_t^{N,n}
  \right\|_1
  \le
  2d_{\mathrm{TV}}
  \left(
    \mathbf P_{N,t}^{X^{1:n}\mid Y^{1:N}},
    \bigotimes_{j=1}^n
    \mathbf q_t^{X\mid Y=Y^j}
  \right)\le
  \left[
    2
    \Ent\left(
      \mathbf P_{N,t}^{X^{1:n}\mid Y^{1:N}}
      \mathrel{\bigg|}
      \bigotimes_{j=1}^n
      \mathbf q_t^{X\mid Y=Y^j}
    \right)
  \right]^{1/2}.
\end{aligned}
\]
Indeed, the evaluation map
\[
  (x^1,\ldots,x^n)
  \longmapsto
  x_t^1\otimes\cdots\otimes x_t^n
\]
takes values in density matrices and therefore has trace norm one.  Taking 
expectations, using Jensen's inequality, and then applying
\eqref{eq:eta-block-conditional-entropy} and
Proposition~\ref{prop:eta-entropy} proves the first bound in
\eqref{eq:eta-conditional-factorization-bound}.  The expectations under
\(\Q^{N,\eta}\) and \(\Q^{N,\mathrm{full}}\) agree because the random
variables are \(\cY_t^N\)-measurable.  Finally,
\(N/r_N\) remains bounded for fixed \(n\), and
Proposition~\ref{prop:eta-full-chaos} gives the stated rate.
\end{proof}

Combining \eqref{eq:eta-first-term-bound} with
Proposition~\ref{prop:eta-conditional-factorization}, we obtain
\begin{equation}
  \sup_{0\le t\le T}
  \E_{\Q^{N,\eta}}
  \left[
    \left\|
      \rho_t^{N:n,\eta}-\overline\Gamma_t^{N,n}
    \right\|_1
  \right]
  \le
  C_{\eta,T,n}N^{-1/8}.
  \label{eq:eta-intermediate-chaos}
\end{equation}

%%%%%%%%%%%%%%%%%%%%%%%%%%%%%%%%%%%%%%%%%%%%%%%%%%%%%%%%%%%%%%%%%%%%%%%%%%%%%%%
\subsection{Identification of the limiting filter and conclusion}
\label{subsec:eta-return-to-target}
%%%%%%%%%%%%%%%%%%%%%%%%%%%%%%%%%%%%%%%%%%%%%%%%%%%%%%%%%%%%%%%%%%%%%%%%%%%%%%%

We now select a canonical version of the conditional barycenter.  Let
\(\mathbb W_T^m\) denote Wiener measure on \(\mathsf Y_T\), and let
\(Y=(Y^1,\ldots,Y^m)\) be the coordinate process.  On the canonical Wiener
space, consider the linear equation
\begin{equation}
\begin{aligned}
  \d\widehat{\overline\gamma}_t
  =
  \mathfrak l_t(\widehat{\overline\gamma}_t)\,\d t
  +
  \sqrt{\eta}
  \sum_{k=1}^m
  \cC_{L^{(k)}}(\widehat{\overline\gamma}_t)\,\d Y_t^k,\qquad 
  \widehat{\overline\gamma}_0
  =
  \gamma_0,
\end{aligned}
  \label{eq:eta-local-linear-SDE}
\end{equation}
and set
\begin{equation}
  \Pi_t(Y)
  :=
  \frac{
    \widehat{\overline\gamma}_t
  }{
    \Tr(\widehat{\overline\gamma}_t)
  }.
  \label{eq:eta-Pi-def}
\end{equation}
Pathwise uniqueness gives a progressively measurable solution map
\(Y\longmapsto\Pi(Y)\), defined \(\mathbb W_T^m\)-almost surely.  We fix an
arbitrary Borel extension to all of \(\mathsf Y_T\).



\begin{lemma}
\label{lem:eta-barycenter-identification}
Let \(\mathbf q_T^Y\) denote the \(\mathsf Y_T\)-marginal of
\(\mathbf q_T\). Then \(\mathbf q_T^Y\) is equivalent to
\(\mathbb W_T^m\), and the version in
\eqref{eq:eta-bar-gamma-barycenter} may be chosen so that
\begin{equation}
  \overline\gamma_t(Y)
  =
  \Pi_t(Y)
  =
  \E_{\mathbf q_T}
  \left[
    X_t\mid\sigma(Y_s:0\le s\le t)
  \right].
  \label{eq:eta-barycenter-identification}
\end{equation}
Moreover, under Wiener measure, and hence under every equivalent measure,
\(\Pi\) satisfies
\begin{equation}
\begin{aligned}
  \d\Pi_t
  =
  \mathfrak l_t(\Pi_t)\,\d t+
  \sqrt{\eta}
  \sum_{k=1}^m
  \cE_{L^{(k)}}(\Pi_t)
  \left(
    \d Y_t^k
    -
    \sqrt{\eta}\,b^{(k)}(\Pi_t)\,\d t
  \right).
\end{aligned}
  \label{eq:eta-local-observation-SDE}
\end{equation}
\end{lemma}

\begin{proof}
Apply the one-particle version of the construction in
Proposition~\ref{prop:eta-conditional-reduction}.  More explicitly, on a
reference space carrying independent Brownian motions \(Y^k,U^k\), let
\(\widehat R\) solve the full linear equation obtained from
\eqref{eq:eta-full-linear-N} by replacing \(\mathfrak L_N\) with
\(\mathfrak l_t\).  Set
\[
  M_t:=\Tr(\widehat R_t),
  \qquad
  R_t:=\frac{\widehat R_t}{M_t},
  \qquad
  \frac{\d\Q^{\mathrm{full}}}{\d\Q}\bigg|_{\cF_T}:=M_T .
\]
This is the normalization and change of measure by the trace.  As in
Proposition~\ref{prop:eta-conditional-reduction}, \(M\) is a strictly positive
martingale, and under \(\Q^{\mathrm{full}}\) the pair \((R,Y)\) has law
\(\mathbf q_T\).

Let
\(
  \cY_t:=\sigma(Y_s:0\le s\le t).
\)
Projecting the full linear equation onto \(\cY_t\) gives
\(
  \E_{\Q}[\widehat R_t\mid\cY_t]
  =
  \widehat{\overline\gamma}_t,
\)
where \(\widehat{\overline\gamma}\) is the solution of
\eqref{eq:eta-local-linear-SDE}.  Bayes' formula then yields
\[
\begin{aligned}
  \E_{\mathbf q_T}[X_t\mid\cY_t]
  &=
  \E_{\Q^{\mathrm{full}}}[R_t\mid\cY_t]  =
  \frac{
    \E_{\Q}[M_T R_t\mid\cY_t]
  }{
    \E_{\Q}[M_T\mid\cY_t]
  }
  =
  \frac{
    \E_{\Q}[\widehat R_t\mid\cY_t]
  }{
    \Tr\left(\E_{\Q}[\widehat R_t\mid\cY_t]\right)
  }
  =
  \frac{\widehat{\overline\gamma}_t}
       {\Tr(\widehat{\overline\gamma}_t)}
  =
  \Pi_t(Y).
\end{aligned}
\]
This proves \eqref{eq:eta-barycenter-identification}.  It remains to identify the observation law.  For every bounded measurable
function \(F\) on \(\mathsf Y_T\),
\[
\begin{aligned}
  \int F(y)\,\mathbf q_T^Y(\d y)
  &=
  \E_{\Q^{\mathrm{full}}}[F(Y)]
  =
  \E_{\Q}[F(Y)M_T]  =
  \E_{\Q}
  \left[
    F(Y)\E_{\Q}[M_T\mid\cY_T]
  \right]
  =
  \E_{\Q}
  \left[
    F(Y)\Tr(\widehat{\overline\gamma}_T)
  \right].
\end{aligned}
\]
Since \(Y\) has law \(\mathbb W_T^m\) under \(\Q\), we obtain
\[
  \frac{\d\mathbf q_T^Y}{\d\mathbb W_T^m}(Y)
  =
  \Tr(\widehat{\overline\gamma}_T).
\]
The density is strictly positive, because it is the conditional expectation of
the strictly positive random variable \(M_T\).  Hence
\(\mathbf q_T^Y\sim\mathbb W_T^m\).  Finally, applying It\^o's formula to
\(
  \Pi_t
\)
in \eqref{eq:eta-local-linear-SDE} gives
\eqref{eq:eta-local-observation-SDE}.
\end{proof}







On \((\Omega,\cY_T^N,\Q^{N,\eta})\), define
\begin{equation}
  \overline\gamma_t^{N,j}
  :=
  \Pi_t(Y^j),
  \qquad
  Y^j=(Y^{1,j},\ldots,Y^{m,j}).
  \label{eq:eta-bar-gamma-N-def}
\end{equation}
Since \(\Q^{N,\eta}\) is equivalent to \(\Q\) on \(\cY_T^N\), each
\(Y^j\)-law is equivalent to Wiener measure.  Hence
Lemma~\ref{lem:eta-barycenter-identification} gives
\begin{equation}
  \overline\Gamma_t^{N,n}
  =
  \overline\gamma_t^{N,1}
  \otimes\cdots\otimes
  \overline\gamma_t^{N,n}.
  \label{eq:eta-bar-Gamma-product}
\end{equation}
Using \eqref{eq:eta-WN-def} in
\eqref{eq:eta-local-observation-SDE}, we obtain, under
\(\Q^{N,\eta}\),
\begin{equation}
\begin{aligned}
  \d\overline\gamma_t^{N,j}
  &=
  \mathfrak l_t(\overline\gamma_t^{N,j})\,\d t
  +
  \sqrt{\eta}
  \sum_{k=1}^m
  \cE_{L^{(k)}}(\overline\gamma_t^{N,j})\,\d W_t^{N,k,j}
  +
  \eta
  \sum_{k=1}^m
  \cE_{L^{(k)}}(\overline\gamma_t^{N,j})
  \left(
    b_j^{N,k}(\rho_t^{N,\eta})
    -
    b^{(k)}(\overline\gamma_t^{N,j})
  \right)\d t.
\end{aligned}
  \label{eq:eta-bar-gamma-under-Qeta}
\end{equation}
For \(j=1,\ldots,N\), let \(\gamma^{N,\eta,j}\) solve, on the same
probability space and with the same Brownian motions,
\begin{equation}
\begin{aligned}
  \d\gamma_t^{N,\eta,j}
  =
  \mathfrak l_t(\gamma_t^{N,\eta,j})\,\d t
  +
  \sqrt{\eta}
  \sum_{k=1}^m
  \cE_{L^{(k)}}(\gamma_t^{N,\eta,j})\,\d W_t^{N,k,j},\qquad 
  \gamma_0^{N,\eta,j}
  =
  \gamma_0.
\end{aligned}
  \label{eq:eta-MF-target}
\end{equation}
The processes \(\gamma^{N,\eta,1},\ldots,\gamma^{N,\eta,N}\) are independent
and identically distributed under \(\Q^{N,\eta}\).  Taking expectations in
\eqref{eq:eta-MF-target} and using uniqueness for
\eqref{eq:mean-state-equation} gives
\[
  \E_{\Q^{N,\eta}}[\gamma_t^{N,\eta,j}]
  =
  \xi_t,
\]
so \eqref{eq:eta-MF-target} is precisely the limiting equation
\eqref{eq:MF-Belavkin}.

\begin{lemma}
\label{lem:eta-projected-exchangeability}
Under $\Q^{N,\eta}$, the family
\[
  \bigl(
    \rho^{N,\eta},
    (\overline\gamma^{N,j},\gamma^{N,\eta,j},Y^j)_{j=1}^N
  \bigr)
\]
is exchangeable.
\end{lemma}

\begin{proof}
The conditional expectation in
\eqref{eq:eta-normalized-conditional} and the normalization by
$Z^{N,\eta}$ commute with simultaneous permutations of the particle labels.
Moreover, the density $Z_T^{N,\eta}$ is permutation invariant.  Hence
Lemma~\ref{lem:eta-full-exchangeability} implies exchangeability of
$(\rho^{N,\eta},(Y^j)_{j=1}^N)$ under $\Q^{N,\eta}$.  The map
$Y^j\longmapsto\overline\gamma^{N,j}=\Pi(Y^j)$ is the same for every particle,
and the equations \eqref{eq:eta-MF-target} are driven by the correspondingly
permuted innovation processes with identical initial data.  Uniqueness in
law completes the proof.
\end{proof}

\begin{lemma}
\label{lem:eta-local-comparison}
For every \(T>0\), there exists \(C_{\eta,T}<\infty\), independent of
\(N\) and \(j\), such that
\begin{equation}
\begin{aligned}
  &\E_{\Q^{N,\eta}}
  \left[
    \sup_{0\le t\le T}
    \left\|
      \overline\gamma_t^{N,j}
      -
      \gamma_t^{N,\eta,j}
    \right\|_1^2
  \right]\le
  C_{\eta,T}
  \int_0^T
  \E_{\Q^{N,\eta}}
  \left[
    \left\|
      \rho_s^{N:(j),\eta}
      -
      \overline\gamma_s^{N,j}
    \right\|_1^2
  \right]\d s.
\end{aligned}
  \label{eq:eta-local-comparison}
\end{equation}
\end{lemma}

\begin{proof}
Set
\(
  \Delta\gamma_t^j
  :=
  \overline\gamma_t^{N,j}
  -
  \gamma_t^{N,\eta,j}.
\)
Subtracting \eqref{eq:eta-MF-target} from
\eqref{eq:eta-bar-gamma-under-Qeta} gives a linear Lipschitz drift, the
martingale term
\[
  \sqrt{\eta}
  \sum_{k=1}^m
  \left(
    \cE_{L^{(k)}}(\overline\gamma_t^{N,j})
    -
    \cE_{L^{(k)}}(\gamma_t^{N,\eta,j})
  \right)\d W_t^{N,k,j},
\]
and the additional drift
\[
  \eta
  \sum_{k=1}^m
  \cE_{L^{(k)}}(\overline\gamma_t^{N,j})
  \left(
    b_j^{N,k}(\rho_t^{N,\eta})
    -
    b^{(k)}(\overline\gamma_t^{N,j})
  \right).
\]
On the compact state space \(\cS(\HH)\), every map
\(\cE_{L^{(k)}}\) is bounded and Lipschitz.  The
Burkholder--Davis--Gundy inequality in Hilbert--Schmidt norm, followed by
Gronwall's lemma, therefore gives
\[
\begin{aligned}
  &\E_{\Q^{N,\eta}}
  \left[
    \sup_{0\le t\le T}
    \|\Delta\gamma_t^j\|_2^2
  \right] \le
  C_{\eta,T}
  \sum_{k=1}^m
  \int_0^T
  \E_{\Q^{N,\eta}}
  \left|
    b_j^{N,k}(\rho_s^{N,\eta})
    -
    b^{(k)}(\overline\gamma_s^{N,j})
  \right|^2\d s.
\end{aligned}
\]
For every \(k\),  it holds that
\[
\begin{aligned}
  \left|
    b_j^{N,k}(\rho_s^{N,\eta})
    -
    b^{(k)}(\overline\gamma_s^{N,j})
  \right| \le
  \|A^{(k)}\|
  \left\|
    \rho_s^{N:(j),\eta}
    -
    \overline\gamma_s^{N,j}
  \right\|_1.
\end{aligned}
\]
Equivalence between the norms $\|\cdot\|_1$ and $\|\cdot\|_2$ on \(\cL(\HH)\)
proves \eqref{eq:eta-local-comparison}.
\end{proof}

We are now ready to complete the proof of Proposition~\ref{prop:main-exact-eta-less-one}.

\begin{proof}[Proof of Proposition~\ref{prop:main-exact-eta-less-one}]
By \eqref{eq:eta-intermediate-chaos} with \(n=1\), Lemma~\ref{lem:eta-projected-exchangeability}, and the
bound \(\|\rho-\sigma\|_1^2\le2\|\rho-\sigma\|_1\) for density matrices,
\[
\begin{aligned}
  \int_0^T
  \E_{\Q^{N,\eta}}
  \left[
    \left\|
      \rho_s^{N:(j),\eta}
      -
      \overline\gamma_s^{N,j}
    \right\|_1^2
  \right]\d s\le
  C_{\eta,T}N^{-1/8}.
\end{aligned}
\]
Lemma~\ref{lem:eta-local-comparison} and Cauchy--Schwarz therefore yield
\begin{equation}
  \E_{\Q^{N,\eta}}
  \left[
    \sup_{0\le t\le T}
    \left\|
      \overline\gamma_t^{N,j}
      -
      \gamma_t^{N,\eta,j}
    \right\|_1
  \right]
  \le
  C_{\eta,T}N^{-1/16}.
  \label{eq:eta-local-process-rate}
\end{equation}
Set
\(
  \Gamma_t^{N,\eta,n}
  :=
  \gamma_t^{N,\eta,1}
  \otimes\cdots\otimes
  \gamma_t^{N,\eta,n}.
\)
The telescopic inequality for tensor products gives
\[
  \left\|
    \overline\Gamma_t^{N,n}
    -
    \Gamma_t^{N,\eta,n}
  \right\|_1
  \le
  \sum_{j=1}^n
  \left\|
    \overline\gamma_t^{N,j}
    -
    \gamma_t^{N,\eta,j}
  \right\|_1.
\]
Combining this estimate with \eqref{eq:eta-intermediate-chaos} and
\eqref{eq:eta-local-process-rate} proves
\[
  \sup_{0\le t\le T}
  \E_{\Q^{N,\eta}}
  \left[
    \left\|
      \rho_t^{N:n,\eta}
      -
      \Gamma_t^{N,\eta,n}
    \right\|_1
  \right]
  \le
  C_{\eta,T,n}N^{-1/16}.
\]
The pair \((\rho^{N,\eta},\gamma^{N,\eta,1},\ldots,\gamma^{N,\eta,n})\)
is a weak realization of the synchronously coupled equations
\eqref{eq:N-Belavkin} and \eqref{eq:MF-Belavkin}.  Uniqueness in law
therefore identifies its expectation with that in
Definition~\ref{def:propagation-of-chaos}, proving
the desired inequality \eqref{eq:eta-main-rate}.
\end{proof}

\begin{remark}
{\rm (i)} The second, unobserved channel has two roles.  It restores perfect total
observation, so that the result of Section~\ref{sec:purification} applies,
and it supplies the compensating Brownian drift in
\eqref{eq:eta-controls-def}.  The latter changes the drift of the observed
record from the interacting quantity \(b_j^{N,k}(R^N)\) to the one-particle
quantity \(b^{(k)}(X^j)\) without changing the SDE for \(X^j\).

\noindent {\rm (ii)} However,  this analysis yields an entropy constant containing the factor \((1-\eta)^{-1}\),  and their estimates
are not uniform as \(\eta\uparrow1\).  The endpoint \(\eta=1\)
must be treated directly by purification in Section~\ref{sec:purification}.
\end{remark}



\section{Propagation of chaos: stability approach}
\label{sec:initial-perturbation}


We now fix \(0<\eta\le1\) and \(\gamma_0\in\cS(\HH)\), and extend the exact-product results of Propositions~\ref{prop:main-exact-eta-one} and~\ref{prop:main-exact-eta-less-one} to initial states satisfying the strong approximate tensorization condition
\[
  \left\|
    \rho_0^N-\gamma_0^{\otimes N}
  \right\|_1
  \longrightarrow0.
\]
The argument is the same for perfect and imperfect measurement efficiency,  and  the main difficulty is to obtain a stability estimate whose constant is
independent of \(N\).  Indeed, suppose that one directly subtracts two
solutions of the nonlinear \(N\)-particle equation driven by the same
Brownian motions.  A standard Hilbert--Schmidt estimate then contains the
quadratic-variation term
\[
  \eta
  \sum_{k=1}^m\sum_{j=1}^N
  \left\|
    \cE_{L_j^{(k)}}(\rho_t^{N,\alpha})
    -
    \cE_{L_j^{(k)}}(\rho_t^{N,\beta})
  \right\|_2^2.
\]
Estimating the summands separately introduces a factor of order \(N\).
Moreover, the Hilbert--Schmidt norm of a lifted observable may depend on
\(\dim(\HH_N)=\dim(\HH)^N\).  The resulting Gronwall constants therefore
grow with the full many-particle dimension and are useless in the
mean-field limit.

To avoid this problem, we compare the two initial conditions through the
linear Zakai equation.  Its solution map is positive and
preserves the trace in expectation.  These two structural properties yield
an expected trace-norm contraction with constant one, independently of the
number of particles, the number of noise terms and the dimension of
\(\HH_N\).



\subsection{A common linear reference and uniform
\texorpdfstring{\(N\)}{N}-particle stability}
\label{subsec:common-reference}

We use the same coefficients as in Section~\ref{sec:inefficient}.  For \(X\in\cL(\HH_N)\),  recall 
\[
  \mathfrak L_N(X)=-\i[H^N,X]+\sum_{k=1}^m\sum_{j=1}^N\cD_{L_j^{(k)}}(X),
\]
and 
\[
  A^{(k)}=L^{(k)}+(L^{(k)})^*,
  \qquad
  b_j^{N,k}(R)=\Tr(A_j^{(k)}R),
  \qquad R\in\cS(\HH_N).
\]
Fix \(T>0\), and let
\(
  (\Omega,\cF,(\cF_t)_{0\le t\le T},\Q)
\)
carry independent real-valued Brownian motions
\[
  Y^{k,j},
  \qquad
   k=1,\ldots, m,
  \quad
   j=1, \ldots, N.
\]
This is a new reference space, independent of the construction in
Section~\ref{sec:inefficient}.  We take $(\cF_t)_{0\le t\le T}$ to be the usual augmentation of their
natural filtration.
For an initial condition
\(
  \theta_0^N\in\cS(\HH_N),
\)
let \(\widehat\rho^{N,\theta}\) solve
\begin{equation}
\begin{aligned}
  \d\widehat\rho_t^{N,\theta}
  =
  \mathfrak L_N
  \bigl(
    \widehat\rho_t^{N,\theta}
  \bigr)\,\d t+
  \sqrt{\eta}
  \sum_{k=1}^m\sum_{j=1}^N
  \cC_{L_j^{(k)}}
  \bigl(
    \widehat\rho_t^{N,\theta}
  \bigr)\,\d Y_t^{k,j},
  \qquad 
  \widehat\rho_0^{N,\theta}
  =
  \theta_0^N.
\end{aligned}
  \label{eq:reference-linear}
\end{equation}
The superscript \(\theta\) is only a label for the initial condition.  The
coefficients and the reference Brownian motions are the same for all
choices of \(\theta_0^N\).

The next lemma is the one-channel analogue of Lemma~\ref{lem:full-linear-reference}.  It is obtained by projecting the fully observed reference evolution onto the filtration generated by \(Y\).

\begin{lemma}
\label{lem:CP-reference}
For every \(t\in[0,T]\), equation \eqref{eq:reference-linear} defines a
random positive linear map $ \Phi_t^{N,\eta}:\cL(\HH_N)\longrightarrow\cL(\HH_N)$ by
\[
  \Phi_t^{N,\eta}(\theta_0^N)
  :=
  \widehat\rho_t^{N,\theta}.
\]
Moreover, for every \(B\in\cL^+(\HH_N)\),
\begin{equation}
  \E_{\Q}\left[
    \Tr\bigl(\Phi_t^{N,\eta}(B)\bigr)
  \right]
  =
  \Tr(B).
  \label{eq:reference-mean-trace}
\end{equation}
\end{lemma}

\begin{proof}
This is the observed-channel projection of Lemma~\ref{lem:full-linear-reference}.  More precisely, on an auxiliary extension of the present reference space, add independent Brownian motions \(U^{k,j}\), and let \(\widehat R^{N,B}\) solve the full linear equation \eqref{eq:eta-full-linear-N} with initial condition \(B\).  By the projection argument used in Proposition~\ref{prop:eta-conditional-reduction},
\[
  \Phi_t^{N,\eta}(B)
  =
  \widehat\rho_t^{N,B}
  =
  \E_{\Q}\left[
    \widehat R_t^{N,B}\mid\cF_t
  \right].
\]
Lemma~\ref{lem:full-linear-reference} gives complete positivity of \(B\mapsto\widehat R_t^{N,B}\) and preservation of the trace in expectation.  Conditional expectation preserves positivity of block matrices, hence complete positivity, and the tower property gives \eqref{eq:reference-mean-trace}.
\end{proof}

Set
\begin{equation}
  Z_t^{N,\theta}
  :=
  \Tr\bigl(\widehat\rho_t^{N,\theta}\bigr),
  \qquad
  \rho_t^{N,\theta}
  :=
  \frac{
    \widehat\rho_t^{N,\theta}
  }{
    Z_t^{N,\theta}
  }.
  \label{eq:perturbation-normalization}
\end{equation}
The process \(\widehat\rho^{N,\theta}\) is positive by
Lemma~\ref{lem:CP-reference}.  Moreover,
\[
  Z_t^{N,\theta}>0
  \qquad
  \Q\text{-almost surely}.
\]
Indeed, on the auxiliary extension used in the proof of Lemma~\ref{lem:CP-reference}, the representation \eqref{eq:eta-full-propagator} shows that the full unnormalized trace is strictly positive.  Its conditional expectation is therefore strictly positive as well.

Taking the trace in \eqref{eq:reference-linear} yields
\begin{equation}
\begin{aligned}
  \d Z_t^{N,\theta}
  &=
  \sqrt{\eta}\,
  Z_t^{N,\theta}
  \sum_{k=1}^m\sum_{j=1}^N
  b_j^{N,k}(\rho_t^{N,\theta})\,\d Y_t^{k,j}.
\end{aligned}
  \label{eq:perturbation-Z-SDE}
\end{equation}
Since
\[
  \left|
    b_j^{N,k}(R)
  \right|
  \le
  \|A^{(k)}\|,
  \qquad
  R\in\cS(\HH_N),
\]
Novikov's condition holds on every finite time interval.  Thus
\(Z^{N,\theta}\) is a strictly positive martingale with mean one.  Define
the probability measure
\begin{equation}
  \frac{\d\Q^{N,\theta}}{\d\Q}
  \bigg|_{\cF_T}
  :=
  Z_T^{N,\theta}.
  \label{eq:perturbation-Qtheta}
\end{equation}
Under \(\Q^{N,\theta}\), the processes
\begin{equation}
  W_t^{N,\theta,k,j}
  :=
  Y_t^{k,j}
  -
  \sqrt{\eta}
  \int_0^t
  b_j^{N,k}(\rho_s^{N,\theta})\,\d s
  \label{eq:perturbation-Wtheta}
\end{equation}
form a family of independent Brownian motions.  Normalizing
\eqref{eq:reference-linear} gives
\begin{equation}
\begin{aligned}
  \d\rho_t^{N,\theta}
  =
  \mathfrak L_N(\rho_t^{N,\theta})\,\d t+
  \sqrt{\eta}
  \sum_{k=1}^m\sum_{j=1}^N
  \cE_{L_j^{(k)}}(\rho_t^{N,\theta})
  \,\d W_t^{N,\theta,k,j}.
\end{aligned}
  \label{eq:perturbation-nonlinear-N}
\end{equation}
Hence, under \(\Q^{N,\theta}\), the process
\(\rho^{N,\theta}\) is a weak solution of the \(N\)-particle Belavkin
equation with initial condition \(\theta_0^N\).

\begin{lemma}
\label{lem:linear-contraction-final}
Let
\(
  \alpha_0^N,\beta_0^N\in\cS(\HH_N).
\)
Then, for every \(t\in[0,T]\),
\begin{equation}
  \E_{\Q}
  \left[
    \left\|
      \widehat\rho_t^{N,\alpha}
      -
      \widehat\rho_t^{N,\beta}
    \right\|_1
  \right]
  \le
  \left\|
    \alpha_0^N-\beta_0^N
  \right\|_1.
  \label{eq:linear-contraction-final}
\end{equation}
\end{lemma}

\begin{proof}
Let
\[
  \alpha_0^N-\beta_0^N
  =
  D_+-D_-
\]
be the Jordan decomposition, where
\(
  D_+,D_-\ge0\) and \(
  D_+D_-=0.
\)
By linearity,
\(
  \widehat\rho_t^{N,\alpha}
  -
  \widehat\rho_t^{N,\beta}
  =
  \Phi_t^{N,\eta}(D_+)
  -
  \Phi_t^{N,\eta}(D_-).
\)
The positivity of \(\Phi_t^{N,\eta}\) thus gives
\[
\begin{aligned}
  \left\|
    \widehat\rho_t^{N,\alpha}
    -
    \widehat\rho_t^{N,\beta}
  \right\|_1
  &\le
  \Tr\bigl(\Phi_t^{N,\eta}(D_+)\bigr)
  +
  \Tr\bigl(\Phi_t^{N,\eta}(D_-)\bigr).
\end{aligned}
\]
Taking expectations and applying
\eqref{eq:reference-mean-trace}, we obtain
\[
\begin{aligned}
  \E_{\Q}
  \left[
    \left\|
      \widehat\rho_t^{N,\alpha}
      -
      \widehat\rho_t^{N,\beta}
    \right\|_1
  \right]\le
  \Tr(D_+)+\Tr(D_-)=
  \left\|
    \alpha_0^N-\beta_0^N
  \right\|_1.
\end{aligned}
\]
\end{proof}

\begin{lemma}
\label{lem:normalization-final}
Let \(B,C\ge0\) be nonzero, and set
\(
  b:=\Tr(B)\) and \(
  c:=\Tr(C).
\)
Then
\begin{equation}
  b
  \left\|
    \frac{B}{b}-\frac{C}{c}
  \right\|_1
  \le
  2\|B-C\|_1.
  \label{eq:normalization-inequality}
\end{equation}
\end{lemma}

\begin{proof}
We have
\[
\begin{aligned}
  b
  \left\|
    \frac{B}{b}-\frac{C}{c}
  \right\|_1
  =
  \left\|
    B-\frac{b}{c}C
  \right\|_1
  \le
  \|B-C\|_1
  +
  \left|
    1-\frac{b}{c}
  \right|
  \|C\|_1
  =
  \|B-C\|_1+|b-c|.
\end{aligned}
\]
Since
\(
  |b-c|
  \le
  \|B-C\|_1,
\)
the result follows.
\end{proof}

For the remainder of the section, set
\begin{equation}
  d_N
  :=
  \left\|
    \alpha_0^N-\beta_0^N
  \right\|_1.
  \label{eq:perturbation-dN}
\end{equation}
For $1\le n\le N$, we use the notation
\(
  \rho_t^{N:n,\theta}
  :=
  \Tr_{[N]\setminus[n]}(\rho_t^{N,\theta}).
\)

\begin{proposition}
\label{prop:common-reference-final}
For every \(t\in[0,T]\),
\begin{equation}
  \E_{\Q^{N,\alpha}}
  \left[
    \left\|
      \rho_t^{N,\alpha}
      -
      \rho_t^{N,\beta}
    \right\|_1
  \right]
  \le
  2d_N.
  \label{eq:uniform-N-stability}
\end{equation}
The same estimate holds for every marginal:
\begin{equation}
  \E_{\Q^{N,\alpha}}
  \left[
    \left\|
      \rho_t^{N:n,\alpha}
      -
      \rho_t^{N:n,\beta}
    \right\|_1
  \right]
  \le
  2d_N,
  \qquad
  1\le n\le N.
  \label{eq:uniform-marginal-stability}
\end{equation}
\end{proposition}

\begin{proof}
Apply Lemma~\ref{lem:normalization-final} with
\(
  B=\widehat\rho_t^{N,\alpha}\) and \(
  C=\widehat\rho_t^{N,\beta}.
\)
Multiplying by the density \(Z_t^{N,\alpha}\) and taking
\(\Q\)-expectations gives
\[
\begin{aligned}
  \E_{\Q^{N,\alpha}}
  \left[
    \left\|
      \rho_t^{N,\alpha}
      -
      \rho_t^{N,\beta}
    \right\|_1
  \right]
  =
  \E_{\Q}
  \left[
    Z_t^{N,\alpha}
    \left\|
      \rho_t^{N,\alpha}
      -
      \rho_t^{N,\beta}
    \right\|_1
  \right]
  \le
  2
  \E_{\Q}
  \left[
    \left\|
      \widehat\rho_t^{N,\alpha}
      -
      \widehat\rho_t^{N,\beta}
    \right\|_1
  \right]
  \le
  2d_N.
\end{aligned}
\]
The marginal estimate follows from trace-norm contractivity of the partial
trace.
\end{proof}

Under \(\Q^{N,\alpha}\), only \(\rho^{N,\alpha}\) is asserted to solve the
nonlinear equation with initial condition \(\alpha_0^N\).  The process
\(\rho^{N,\beta}\) in
\eqref{eq:uniform-N-stability} is the normalized solution of the same
linear reference equation, evaluated along the same observation paths, but
with initial condition \(\beta_0^N\).  This distinction is the reason for
introducing the changed measures below.

For \(t\le T\), let
\(
  \Q_t^{N,\theta}
  :=
  \Q^{N,\theta}\big|_{\cF_t}.
\)

\begin{proposition}
\label{prop:TV-final}
For every \(t\in[0,T]\),
\begin{equation}
\begin{aligned}
  d_{\mathrm{TV}}
  \left(
    \Q_t^{N,\alpha},
    \Q_t^{N,\beta}
  \right)
  =
  \frac12
  \E_{\Q}
  \left[
    \left|
      Z_t^{N,\alpha}
      -
      Z_t^{N,\beta}
    \right|
  \right]
 \le
  \frac12d_N.
\end{aligned}
  \label{eq:perturbation-TV-bound}
\end{equation}
\end{proposition}

\begin{proof}
Since \(Z^{N,\theta}\) is a martingale, it is the density of
\(\Q_t^{N,\theta}\) with respect to \(\Q|_{\cF_t}\).  The first equality
is therefore the density formula for total variation.  Moreover,
\[
\begin{aligned}
  \left|
    Z_t^{N,\alpha}
    -
    Z_t^{N,\beta}
  \right|
  =
  \left|
    \Tr\left(
      \widehat\rho_t^{N,\alpha}
      -
      \widehat\rho_t^{N,\beta}
    \right)
  \right|\le
  \left\|
    \widehat\rho_t^{N,\alpha}
    -
    \widehat\rho_t^{N,\beta}
  \right\|_1.
\end{aligned}
\]
The conclusion follows from
Lemma~\ref{lem:linear-contraction-final}.
\end{proof}

In particular, if \(F\) is \(\cF_t\)-measurable and
\(
  a\le F\le b,
\)
then
\begin{equation}
  \left|
    \E_{\Q^{N,\alpha}}[F]
    -
    \E_{\Q^{N,\beta}}[F]
  \right|
  \le
  (b-a)
  d_{\mathrm{TV}}
  \left(
    \Q_t^{N,\alpha},
    \Q_t^{N,\beta}
  \right).
  \label{eq:TV-expectation}
\end{equation}

\begin{remark}
The estimates
\eqref{eq:uniform-N-stability} and
\eqref{eq:perturbation-TV-bound} contain no factor depending on
\(\dim(\HH_N)\), \(mN\), or the norm of the \(N\)-particle generator.
This is the essential gain of the linear reference equation.  The argument
uses positivity and trace preservation rather than a Lipschitz estimate for
the nonlinear \(N\)-particle coefficients.
\end{remark}

%%%%%%%%%%%%%%%%%%%%%%%%%%%%%%%%%%%%%%%%%%%%%%%%%%%%%%%%%%%%%%%%%%%%%%%%%%%%%%%
\subsection{Comparison of the limiting processes and proof of
Theorem~\ref{thm:main-approx}}
\label{subsec:limiting-comparison-perturbation}


Recall that \(\xi\) is the deterministic solution of
\eqref{eq:mean-state-equation}, and 
\[
  \mathfrak l_t(X)
  =
  -\i[h+V^{\xi_t},X]
  +
  \sum_{k=1}^m
  \cD_{L^{(k)}}(X)
 \]
 is defined via  \eqref{eq:eta-little-l-def}.  For
\(
  \theta\in\{\alpha,\beta\}
\)
and \(j=1,\ldots,N\), define \(\gamma^{N,\theta,j}\) on the common
reference space by
\begin{equation}
\begin{aligned}
  \d\gamma_t^{N,\theta,j}
  &=
  \mathfrak l_t
  \bigl(
    \gamma_t^{N,\theta,j}
  \bigr)\,\d t+
  \sqrt{\eta}
  \sum_{k=1}^m
  \cE_{L^{(k)}}
  \bigl(
    \gamma_t^{N,\theta,j}
  \bigr)
  \left(
    \d Y_t^{k,j}
    -
    \sqrt{\eta}\,
    b_j^{N,k}
    \bigl(
      \rho_t^{N,\theta}
    \bigr)\,\d t
  \right),
  \\
  \gamma_0^{N,\theta,j}
  &=
  \gamma_0.
\end{aligned}
  \label{eq:perturbation-gamma-observation-form}
\end{equation}
Equivalently, under \(\Q^{N,\theta}\),
\begin{equation}
\begin{aligned}
  \d\gamma_t^{N,\theta,j}
  =
  \mathfrak l_t
  \bigl(
    \gamma_t^{N,\theta,j}
  \bigr)\,\d t+
  \sqrt{\eta}
  \sum_{k=1}^m
  \cE_{L^{(k)}}
  \bigl(
    \gamma_t^{N,\theta,j}
  \bigr)
  \,\d W_t^{N,\theta,k,j}.
\end{aligned}
  \label{eq:perturbation-gamma-physical-form}
\end{equation}
Thus, under \(\Q^{N,\theta}\), the processes
\(
  \gamma^{N,\theta,1},
  \ldots,
  \gamma^{N,\theta,N}
\)
are independent copies of the limiting one-particle equation.  In
particular,
\begin{equation}
  \E_{\Q^{N,\theta}}
  \left[
    \gamma_t^{N,\theta,j}
  \right]
  =
  \xi_t.
  \label{eq:perturbation-gamma-mean}
\end{equation}
For \(1\le n\le N\), set
\begin{equation}
  \Gamma_t^{N,\theta,n}
  :=
  \gamma_t^{N,\theta,1}
  \otimes\cdots\otimes
  \gamma_t^{N,\theta,n}.
  \label{eq:perturbation-Gamma-theta}
\end{equation}
The two processes are related by
\begin{equation}
\begin{aligned}
  \d W_t^{N,\beta,k,j}
  =
  \d W_t^{N,\alpha,k,j}+
  \sqrt{\eta}
  \left(
    b_j^{N,k}(\rho_t^{N,\alpha})
    -
    b_j^{N,k}(\rho_t^{N,\beta})
  \right)\d t.
\end{aligned}
  \label{eq:perturbation-W-relation}
\end{equation}
This identity allows the limiting processes to be compared under the same
measure and with the same Brownian motions.

\begin{lemma}
\label{lem:MF-comparison-final}
For every \(T>0\), there exists a constant \(C_T<\infty\), independent of
\(N\), \(j\) and \(\eta\in(0,1]\), such that
\begin{equation}
  \sup_{0\le t\le T}
  \E_{\Q^{N,\alpha}}
  \left[
    \left\|
      \gamma_t^{N,\alpha,j}
      -
      \gamma_t^{N,\beta,j}
    \right\|_1
  \right]
  \le
  C_Td_N.
  \label{eq:MF-comparison-final}
\end{equation}
\end{lemma}

\begin{proof}
For \(1\le k\le m\), set
\[
  \delta b_t^{N,k,j}:=b_j^{N,k}(\rho_t^{N,\alpha})-b_j^{N,k}(\rho_t^{N,\beta}).
\]
Under \(\Q^{N,\alpha}\), equation \eqref{eq:perturbation-W-relation} gives
\begin{equation}
  \d\gamma_t^{N,\beta,j}=\mathfrak l_t(\gamma_t^{N,\beta,j})\,\d t+\sqrt{\eta}\sum_{k=1}^m\cE_{L^{(k)}}(\gamma_t^{N,\beta,j})\,\d W_t^{N,\alpha,k,j}+\eta\sum_{k=1}^m\cE_{L^{(k)}}(\gamma_t^{N,\beta,j})\delta b_t^{N,k,j}\,\d t.
  \label{eq:perturbation-gamma-beta-under-alpha}
\end{equation}
Let
\[
  \Delta\gamma_t^j:=\gamma_t^{N,\alpha,j}-\gamma_t^{N,\beta,j}.
\]
Subtracting \eqref{eq:perturbation-gamma-beta-under-alpha} from the equation for \(\gamma^{N,\alpha,j}\), we obtain
\begin{equation}
  \d\Delta\gamma_t^j=\mathfrak l_t(\Delta\gamma_t^j)\,\d t+\sqrt{\eta}\sum_{k=1}^m\left(\cE_{L^{(k)}}(\gamma_t^{N,\alpha,j})-\cE_{L^{(k)}}(\gamma_t^{N,\beta,j})\right)\d W_t^{N,\alpha,k,j}-\eta\sum_{k=1}^m\cE_{L^{(k)}}(\gamma_t^{N,\beta,j})\delta b_t^{N,k,j}\,\d t.
  \label{eq:perturbation-D-SDE}
\end{equation}
Since \(\xi_t\in\cS(\HH)\), the linear maps \(\mathfrak l_t\) are uniformly bounded on \([0,T]\): there exists \(C_{\mathfrak l}<\infty\) such that \(\|\mathfrak l_t(D)\|_2\le C_{\mathfrak l}\|D\|_2\).  Moreover, the maps \(\cE_{L^{(k)}}\) are bounded and Lipschitz on the compact state space \(\cS(\HH)\).  Thus there exist constants \(C_{\mathrm{Lip}},C_{\mathrm{bd}}<\infty\), depending only on the one-particle coefficients and on \(\dim(\HH)\), such that
\begin{equation}
  \sum_{k=1}^m\left\|\cE_{L^{(k)}}(x)-\cE_{L^{(k)}}(y)\right\|_2^2\le C_{\mathrm{Lip}}\|x-y\|_2^2,\qquad x,y\in\cS(\HH).
  \label{eq:perturbation-E-Lipschitz}
\end{equation}
and
\begin{equation}
  \sum_{k=1}^m\left\|\cE_{L^{(k)}}(x)\right\|_2\le C_{\mathrm{bd}},\qquad x\in\cS(\HH).
  \label{eq:perturbation-E-bounded}
\end{equation}
For \(\varepsilon>0\), define \(f_\varepsilon(D):=\sqrt{\|D\|_2^2+\varepsilon^2}\).  We regard the Hermitian matrices as a finite-dimensional real Hilbert space with the Hilbert--Schmidt inner product.  Applying It\^o's formula to \(f_\varepsilon(\Delta\gamma_t^j)\), using \eqref{eq:perturbation-E-Lipschitz} and \eqref{eq:perturbation-E-bounded}, and using \(0<\eta\le1\), gives
\begin{equation}
  \E_{\Q^{N,\alpha}}\left[f_\varepsilon(\Delta\gamma_t^j)\right]\le\varepsilon+C\int_0^t\E_{\Q^{N,\alpha}}\left[f_\varepsilon(\Delta\gamma_s^j)\right]\d s+C\sum_{k=1}^m\int_0^t\E_{\Q^{N,\alpha}}\left[|\delta b_s^{N,k,j}|\right]\d s,
  \label{eq:perturbation-regularized-estimate}
\end{equation}
where \(C\) is independent of \(N\), \(j\), \(\eta\) and \(\varepsilon\).  By the definition of \(b_j^{N,k}\),
\[
  |\delta b_s^{N,k,j}|=\left|\Tr\left(A^{(k)}\left(\rho_s^{N:(j),\alpha}-\rho_s^{N:(j),\beta}\right)\right)\right|\le\|A^{(k)}\|\left\|\rho_s^{N:(j),\alpha}-\rho_s^{N:(j),\beta}\right\|_1.
\]
By Proposition~\ref{prop:common-reference-final} and contractivity of the partial trace,
\begin{equation}
  \E_{\Q^{N,\alpha}}\left[|\delta b_s^{N,k,j}|\right]\le2\|A^{(k)}\|d_N.
  \label{eq:perturbation-deltab-bound}
\end{equation}
Combining \eqref{eq:perturbation-regularized-estimate} and \eqref{eq:perturbation-deltab-bound}, and then applying Gronwall's lemma, yields
\[
  \sup_{0\le t\le T}\E_{\Q^{N,\alpha}}\left[f_\varepsilon(\Delta\gamma_t^j)\right]\le C_T(\varepsilon+d_N),
\]
where \(C_T\) is independent of \(N\), \(j\) and \(\eta\).  Letting \(\varepsilon\downarrow0\), we obtain
\[
  \sup_{0\le t\le T}\E_{\Q^{N,\alpha}}\left[\|\Delta\gamma_t^j\|_2\right]\le C_Td_N.
\]
Since all norms are equivalent on the finite-dimensional space \(\cL(\HH)\), this proves \eqref{eq:MF-comparison-final}.
\end{proof}

  
  
  
  

We can now formulate the transfer estimate which is the main result of this
section.

\begin{proposition}
\label{prop:uniform-transfer-final}
For every fixed \(n\ge1\) and \(T>0\), there exists
\(C_{T,n}<\infty\), independent of \(N\) and
\(\eta\in(0,1]\), such that
\begin{equation}
\begin{aligned}
  &\sup_{0\le t\le T}
  \E_{\Q^{N,\alpha}}
  \left[
    \left\|
      \rho_t^{N:n,\alpha}
      -
      \Gamma_t^{N,\alpha,n}
    \right\|_1
  \right] \le
  \sup_{0\le t\le T}
  \E_{\Q^{N,\beta}}
  \left[
    \left\|
      \rho_t^{N:n,\beta}
      -
      \Gamma_t^{N,\beta,n}
    \right\|_1
  \right]
  +
  C_{T,n}d_N.
\end{aligned}
  \label{eq:uniform-transfer-final}
\end{equation}
\end{proposition}

\begin{proof}
For every \(t\in[0,T]\),
\begin{align}
  \left\|
    \rho_t^{N:n,\alpha}
    -
    \Gamma_t^{N,\alpha,n}
  \right\|_1
  \le
  \left\|
    \rho_t^{N:n,\alpha}
    -
    \rho_t^{N:n,\beta}
  \right\|_1
  +
  \left\|
    \rho_t^{N:n,\beta}
    -
    \Gamma_t^{N,\beta,n}
  \right\|_1+
  \left\|
    \Gamma_t^{N,\beta,n}
    -
    \Gamma_t^{N,\alpha,n}
  \right\|_1.
  \label{eq:perturbation-three-terms}
\end{align}
By Proposition~\ref{prop:common-reference-final}, the first term satisfies
\begin{equation}
  \E_{\Q^{N,\alpha}}
  \left[
    \left\|
      \rho_t^{N:n,\alpha}
      -
      \rho_t^{N:n,\beta}
    \right\|_1
  \right]
  \le
  2d_N.
  \label{eq:perturbation-first-term}
\end{equation}
For the second term, set
\[
  F_t^{N,\beta}
  :=
  \left\|
    \rho_t^{N:n,\beta}
    -
    \Gamma_t^{N,\beta,n}
  \right\|_1.
\]
Since both arguments are density matrices,
\(
  0\le F_t^{N,\beta}\le2.
\)
Equations
\eqref{eq:TV-expectation} and
\eqref{eq:perturbation-TV-bound} therefore imply
\begin{equation}
\begin{aligned}
  \E_{\Q^{N,\alpha}}
  \left[
    F_t^{N,\beta}
  \right]
  &\le
  \E_{\Q^{N,\beta}}
  \left[
    F_t^{N,\beta}
  \right]
  +
  d_N.
\end{aligned}
  \label{eq:perturbation-second-term}
\end{equation}
Finally, the telescopic inequality for tensor products gives
\begin{equation}
\begin{aligned}
  \left\|
    \Gamma_t^{N,\beta,n}
    -
    \Gamma_t^{N,\alpha,n}
  \right\|_1
  &\le
  \sum_{j=1}^n
  \left\|
    \gamma_t^{N,\beta,j}
    -
    \gamma_t^{N,\alpha,j}
  \right\|_1.
\end{aligned}
  \label{eq:perturbation-tensor-telescopic}
\end{equation}
Lemma~\ref{lem:MF-comparison-final} yields
\begin{equation}
  \sup_{0\le t\le T}
  \E_{\Q^{N,\alpha}}
  \left[
    \left\|
      \Gamma_t^{N,\beta,n}
      -
      \Gamma_t^{N,\alpha,n}
    \right\|_1
  \right]
  \le
  nC_Td_N.
  \label{eq:perturbation-third-term}
\end{equation}
Combining
\eqref{eq:perturbation-first-term},
\eqref{eq:perturbation-second-term} and
\eqref{eq:perturbation-third-term} proves
\eqref{eq:uniform-transfer-final}.
\end{proof}

Finally,  we turn to prove the first main result.

\begin{proof}[Proof of Theorem~\ref{thm:main-approx}]
Apply Proposition~\ref{prop:uniform-transfer-final} with
\(
  \alpha_0^N:=\rho_0^N\) and 
  \(
  \beta_0^N:=\gamma_0^{\otimes N}.
\)
Then
\[
  d_N
  =
  \left\|
    \rho_0^N-\gamma_0^{\otimes N}
  \right\|_1
  =
  \delta_N.
\]
Under \(\Q^{N,\beta}\), the process
\(\rho^{N,\beta}\) has exact tensor-product initial condition.  Therefore
  Propositions~\ref{prop:main-exact-eta-one} and \ref{prop:main-exact-eta-less-one}  give
\[
\begin{aligned}
  &\sup_{0\le t\le T}
  \E_{\Q^{N,\beta}}
  \left[
    \left\|
      \rho_t^{N:n,\beta}
      -
      \Gamma_t^{N,\beta,n}
    \right\|_1
  \right]\le  \frac{C}{N^{1/4}}{\mathbf 1}_{\{\eta=1\}} + \frac{C}{N^{1/16}}{\mathbf 1}_{\{0<\eta<1\}},
\end{aligned}
\]
for some constant $C>0$.  Consequently,
\[
\begin{aligned}
  &\sup_{0\le t\le T}
  \E_{\Q^{N,\alpha}}
  \left[
    \left\|
      \rho_t^{N:n,\alpha}
      -
      \Gamma_t^{N,\alpha,n}
    \right\|_1
  \right]
  \le
   \frac{C}{N^{1/4}}{\mathbf 1}_{\{\eta=1\}} + \frac{C}{N^{1/16}}{\mathbf 1}_{\{0<\eta<1\}}
  +
  C\delta_N.
\end{aligned}
\]
Under \(\Q^{N,\alpha}\), the pair
\[
  \left(
    \rho^{N,\alpha},
    \gamma^{N,\alpha,1},
    \ldots,
    \gamma^{N,\alpha,n}
  \right)
\]
is a weak realization of the synchronously coupled equations
\eqref{eq:N-Belavkin} and \eqref{eq:MF-Belavkin}.  Uniqueness in law
therefore identifies this expectation with the one appearing in
Definition~\ref{def:propagation-of-chaos}.  This fulfills the proof.
\end{proof}




\section[Skew-adjoint measurements: BBGKY hierarchy]{Propagation of chaos for skew-adjoint measurements: stochastic BBGKY hierarchy}
\label{sec:skew-bbgky}





Next we prove Theorem~\ref{thm:main-anti},  and we assume,  throughout this section,  
\[
  (L^{(k)})^*=-L^{(k)},
  \qquad k=1,\ldots,m,
\]
and therefore the $N$-particle   equation  \eqref{eq:N-Belavkin} and the mean-field equation  \eqref{eq:MF-Belavkin} read as 
\begin{equation}
\begin{aligned}
  \d\rho_t^N
  &=
  -\i[H^N,\rho_t^N]\,\d t
  +
  \sum_{k=1}^m\sum_{j=1}^N
  \cD_{L_j^{(k)}}(\rho_t^N)\,\d t
 +
  \sqrt{\eta}
  \sum_{k=1}^m\sum_{j=1}^N
 [L_j^{(k)},  \rho_t^N] \,\d W_t^{k,j}, \\
   \d\gamma_t^j
 & =
  \left(
    -\i[h+V^{\xi_t},\gamma_t^j]
    +
    \sum_{k=1}^m
    \cD_{L^{(k)}}(\gamma_t^j)
  \right)\d t+
  \sqrt{\eta}
  \sum_{k=1}^m
  [L^{(k)},  \gamma_t^j]\,\d W_t^{k,j},
\end{aligned}
 \end{equation}
This linearity is the only special structure used below.

For \(X\in\cL(\HH_{n})\), set
\begin{equation}
 \mathfrak A_n(X):=-\i\left[\sum_{j=1}^n h_j,X\right]+\sum_{k=1}^m\sum_{j=1}^n\cD_{L_j^{(k)}}(X),
  \qquad
  \mathcal M_{k,j}^n(X):=[L_j^{(k)},X].
  \label{eq:skew-local-generator}
\end{equation}
For \(X\in\cL(\HH_{n+1})\), define
\begin{equation}
  \mathcal B_j^n(X):=\Tr_{\{n+1\}}\bigl([V_{j,n+1},X]\bigr),
  \qquad j=1,\ldots,n,
  \label{eq:skew-B-def}
\end{equation}
where \(\Tr_{\{n+1\}}\) denotes the partial trace over the last tensor factor. By duality, for every \(X\in\cL(\HH^{\otimes(n+1)})\),
\begin{equation}
  \left\|
    \mathcal B_j^n(X)
  \right\|_1
  \le
  2\|V\|\,\|X\|_1.
  \label{eq:skew-B-bound}
\end{equation}
Indeed, if \(\|A\|\le1\), then
\[
  \left|\Tr\left(A\,\mathcal B_j^n(X)\right)\right|
  =
  \left|\Tr\left([A\otimes I_{\HH},V_{j,n+1}]X\right)\right|
  \le
  2\|V\|\,\|X\|_1.
\]
For \(n<N\), define the averaged \((n+1)\)-particle marginal
\begin{equation}
  \overline\rho_t^{N:n+1}
  :=
  \frac1{N-n}
  \sum_{\ell=n+1}^N
  \rho_t^{N:[n]\cup\{\ell\}},
  \label{eq:skew-averaged-marginal}
\end{equation}
where each marginal is canonically relabelled as an operator on \(\HH_{n+1}\).

\begin{lemma}
\label{lem:skew-marginal-hierarchy}
For \(1\le n<N\), the marginal \(\rho^{N:n}\) satisfies
\begin{equation}
\begin{aligned}
  \d\rho_t^{N:n}
  &=
  \mathfrak A_n(\rho_t^{N:n})\,\d t
  +
  \sqrt{\eta}
  \sum_{k=1}^m\sum_{j=1}^n
  \mathcal M_{k,j}^n(\rho_t^{N:n})\,\d W_t^{k,j}
  \\
  &\quad
  -
  \i\frac{N-n}{N}
  \sum_{j=1}^n
  \mathcal B_j^n(\overline\rho_t^{N:n+1})\,\d t
  -
  \frac{\i}{N}
  \sum_{1\le i<j\le n}
  [V_{ij},\rho_t^{N:n}]\,\d t.
\end{aligned}
  \label{eq:skew-marginal-hierarchy}
\end{equation}
Moreover, the tensor product \(\Gamma_t^n=\gamma_t^1\otimes\cdots\otimes\gamma_t^n\) satisfies
\begin{equation}
\begin{aligned}
  \d\Gamma_t^n
  &=
 \mathfrak A_n(\Gamma_t^n)\,\d t
  +
  \sqrt{\eta}
  \sum_{k=1}^m\sum_{j=1}^n
  \mathcal M_{k,j}^n(\Gamma_t^n)\,\d W_t^{k,j}
  -
  \i
  \sum_{j=1}^n
  \mathcal B_j^n(\Gamma_t^n\otimes\xi_t)\,\d t.
\end{aligned}
  \label{eq:skew-limit-hierarchy}
\end{equation}
\end{lemma}

\begin{proof}
Take the partial trace of \eqref{eq:N-Belavkin} over the variables \(n+1,\ldots,N\). The terms acting on the first \(n\) particles give the local part in \eqref{eq:skew-marginal-hierarchy}. If \(\ell>n\), then
\[
  \Tr_{\{ \ell\}}\bigl(\cD_{L_\ell^{(k)}}(X)\bigr)=0,
  \qquad
  \Tr_{\{ \ell\}}\bigl([L_\ell^{(k)},X]\bigr)=0,
\]
by cyclicity of the partial trace. Hence all exterior measurement noises disappear. The interactions with both indices larger than \(n\) also vanish after the partial trace. The interactions with both indices at most \(n\) give the last term in \eqref{eq:skew-marginal-hierarchy}, and the \(n(N-n)\) interactions between one retained and one traced particle give the averaged term involving \(\overline\rho_t^{N:n+1}\). This proves \eqref{eq:skew-marginal-hierarchy}. Formula \eqref{eq:skew-limit-hierarchy} follows by applying It\^o's formula to \(\gamma_t^1\otimes\cdots\otimes\gamma_t^n\), using independence of the Brownian motions with different particle labels and the identity \(\mathcal B_j^n(\Gamma_t^n\otimes\xi_t)=[V_j^{\xi_t},\Gamma_t^n]\).
\end{proof}

Let
\[
  \Delta_t^{N,n}:=\rho_t^{N:n}-\Gamma_t^n,
  \qquad
  F_n^N(t):=\E\left[\|\Delta_t^{N,n}\|_1\right].
\]
Subtracting \eqref{eq:skew-limit-hierarchy} from \eqref{eq:skew-marginal-hierarchy} gives
\begin{equation}
\begin{aligned}
  \d\Delta_t^{N,n}
  &=
 \mathfrak A_n(\Delta_t^{N,n})\,\d t
  +
  \sqrt{\eta}
  \sum_{k=1}^m\sum_{j=1}^n
  \mathcal M_{k,j}^n(\Delta_t^{N,n})\,\d W_t^{k,j}
  -
  \i
  \sum_{j=1}^n
  \mathcal B_j^n\left(
    \overline\rho_t^{N:n+1}
    -
    \Gamma_t^n\otimes\xi_t
  \right)\d t
  +
  \mathcal R_t^{N,n}\,\d t,
\end{aligned}
  \label{eq:skew-delta-equation}
\end{equation}
where
\begin{equation}
  \mathcal R_t^{N,n}
  :=
  \frac{\i n}{N}
  \sum_{j=1}^n
  \mathcal B_j^n(\overline\rho_t^{N:n+1})
  -
  \frac{\i}{N}
  \sum_{1\le i<j\le n}
  [V_{ij},\rho_t^{N:n}].
  \label{eq:skew-remainder}
\end{equation}
By \eqref{eq:skew-B-bound} and \(\|\overline\rho_t^{N:n+1}\|_1=\|\rho_t^{N:n}\|_1=1\),
\begin{equation}
  \|\mathcal R_t^{N,n}\|_1
  \le
  \frac{Cn^2}{N},
  \label{eq:skew-remainder-bound}
\end{equation}
where \(C\) depends only on \(\|V\|\).



\begin{proof}[Proof of Theorem~\ref{thm:main-anti}]
Let \(\Phi_{t,s}^n\) be the random solution map of the homogeneous local equation
\[
  \d X_t
  =
  \mathfrak A_n(X_t)\,\d t
  +
  \sqrt{\eta}
  \sum_{k=1}^m\sum_{j=1}^n
  \mathcal M_{k,j}^n(X_t)\,\d W_t^{k,j},
  \qquad X_s=X.
\]
For every \(s\le t\), the map \(\Phi_{t,s}^n\) is completely positive and trace preserving. Indeed, this is the finite-dimensional linear reference equation of Section~\ref{sec:initial-perturbation} restricted to \(n\) particles.  In the present skew-adjoint case the trace process is constant pathwise because \(L^{(k)}+(L^{(k)})^*=0\). Consequently, for every Hermitian \(X\),
\begin{equation}
  \|\Phi_{t,s}^n(X)\|_1
  \le
  \|X\|_1.
  \label{eq:skew-local-contraction}
\end{equation}
The variation-of-constants formula applied to \eqref{eq:skew-delta-equation}, together with \eqref{eq:skew-local-contraction}, yields
\begin{equation}
\begin{aligned}
  F_n^N(t)
  &\le
  \delta_N^n
  +
  Cn
  \int_0^t
  \E\left[
    \left\|
      \overline\rho_s^{N:n+1}
      -
      \Gamma_s^n\otimes\xi_s
    \right\|_1
  \right]\d s
  +
  C\frac{n^2t}{N}.
\end{aligned}
  \label{eq:skew-before-average}
\end{equation}
It remains to estimate the averaged marginal. For \(\ell>n\), write
\[
  \Gamma_t^{n,\ell}
  :=
  \gamma_t^1\otimes\cdots\otimes\gamma_t^n\otimes\gamma_t^\ell,
\]
again relabelled as an operator on \(\HH^{\otimes(n+1)}\). Then
\begin{equation}
\begin{aligned}
  \overline\rho_t^{N:n+1}-\Gamma_t^n\otimes\xi_t
  &=
  \frac1{N-n}
  \sum_{\ell=n+1}^N
  \left(
    \rho_t^{N:[n]\cup\{\ell\}}
    -
    \Gamma_t^{n,\ell}
  \right)
  +
  \Gamma_t^n\otimes
  \left(
    \frac1{N-n}
    \sum_{\ell=n+1}^N
    \gamma_t^\ell
    -
    \xi_t
  \right).
\end{aligned}
  \label{eq:skew-average-decomposition}
\end{equation}
By exchangeability of the coupled family \((\rho^N,\gamma^1,\ldots,\gamma^N)\),
\[
  \E\left[
    \left\|
      \rho_t^{N:[n]\cup\{\ell\}}
      -
      \Gamma_t^{n,\ell}
    \right\|_1
  \right]
  =
  F_{n+1}^N(t).
\]
Moreover, since \(\Gamma_t^n\) has trace norm one,
\[
  \left\|
    \Gamma_t^n\otimes
    \left(
      \frac1{N-n}
      \sum_{\ell=n+1}^N
      \gamma_t^\ell
      -
      \xi_t
    \right)
  \right\|_1
  =
  \left\|
      \frac1{N-n}
      \sum_{\ell=n+1}^N
      \gamma_t^\ell
      -
      \xi_t
  \right\|_1.
\]
The variables \(\gamma_t^\ell\), \(\ell>n\), are i.i.d. with mean \(\xi_t\). Since \(\HH\) is finite-dimensional and \(\|\gamma_t^\ell\|_1=1\),
\begin{equation}
  \sup_{0\le t\le T}
  \E\left[
    \left\|
      \frac1{N-n}
      \sum_{\ell=n+1}^N
      \gamma_t^\ell
      -
      \xi_t
    \right\|_1
  \right]
  \le
  \frac{C}{\sqrt{N-n}}.
  \label{eq:skew-limit-empirical}
\end{equation}
For instance, this follows first in Hilbert--Schmidt norm by independence and then by norm equivalence on \(\cL(\HH)\). Combining \eqref{eq:skew-before-average}--\eqref{eq:skew-limit-empirical}, we obtain, for every fixed \(n<N\),
\begin{equation}
  F_n^N(t)
  \le
  \delta_N^n
  +
  Cn
  \int_0^t
  F_{n+1}^N(s)\,\d s
  +
  C_{T,n}
  \left(
    \frac1{\sqrt{N-n}}
    +
    \frac1N
  \right),
  \qquad 0\le t\le T.
  \label{eq:skew-hierarchy-ineq}
\end{equation}
 We now close the hierarchy by following the iteration argument  in the proof of Theorem~5.2 of Bardos--Golse--Mauser \cite{bardos2000}. Up to increasing the constant, \eqref{eq:skew-hierarchy-ineq} can be written as
\begin{equation}
  F_n^N(t)\le \delta_N^n+Cn\int_0^tF_{n+1}^N(s)\,\d s+\varepsilon_{N,n},
  \qquad 1\le n<N,
  \label{eq:skew-recursive-BGM}
\end{equation}
where
\[
  \varepsilon_{N,n}:=C_T\left(\frac{n}{\sqrt{N-n}}+\frac{n^2}{N}\right).
\]
Iterating \eqref{eq:skew-recursive-BGM}, we obtain the following explicit estimate: for every integer \(K\ge0\) such that \(n+K<N\), and every \(0\le t\le T\),
\begin{equation}
\begin{aligned}
  F_n^N(t)
  &\le
  \sum_{r=0}^K
  \binom{n+r-1}{r}
  (Ct)^r
  \left(
    \delta_N^{n+r}
    +
    \varepsilon_{N,n+r}
  \right)
  +
  2\binom{n+K}{K+1}(Ct)^{K+1}.
\end{aligned}
  \label{eq:skew-BGM-explicit}
\end{equation}
Indeed, after \(r\) iterations the coefficient is
\[
  \frac{n(n+1)\cdots(n+r-1)}{r!}
  =
  \binom{n+r-1}{r},
\]
and the remainder after \(K+1\) iterations is bounded by the trivial estimate \(F_{n+K+1}^N(t)\le2\).

Choose \(\tau>0\) such that \(C\tau<1\). For fixed \(n\) and \(K\), letting \(N\to\infty\) in \eqref{eq:skew-BGM-explicit} gives
\[
  \limsup_{N\to\infty}
  \sup_{0\le t\le\tau}
  F_n^N(t)
  \le
  2\binom{n+K}{K+1}(C\tau)^{K+1},
\]
because \(\delta_N^{n+r}\to0\) and \(\varepsilon_{N,n+r}\to0\) for every fixed \(r\). Since
\[
  \binom{n+K}{K+1}(C\tau)^{K+1}\longrightarrow0
  \qquad\mbox{as }K\to\infty,
\]
we obtain
\[
  \lim_{N\to\infty}
  \sup_{0\le t\le\tau}
  F_n^N(t)=0.
\]
It remains to pass from the short time interval \([0,\tau]\) to an arbitrary finite interval \([0,T]\). We first record the shifted form of the previous estimate. Let \(a\ge0\) and \(a+\tau\le T\). Repeating the variation-of-constants argument on \([a,t]\), instead of on \([0,t]\), gives, for \(a\le t\le a+\tau\),
\begin{equation}
  F_n^N(t)\le F_n^N(a)+Cn\int_a^tF_{n+1}^N(s)\,\d s+\varepsilon_{N,n},
  \qquad
  \varepsilon_{N,n}:=C_{T,n}\left(\frac1{\sqrt{N-n}}+\frac1N\right).
  \label{eq:skew-shifted-recursion}
\end{equation}
Iterating \eqref{eq:skew-shifted-recursion} exactly as before yields, for every \(K\ge0\) such that \(n+K<N\),
\begin{equation}
\begin{aligned}
  \sup_{a\le t\le a+\tau}F_n^N(t)
  &\le
  \sum_{r=0}^K
  \binom{n+r-1}{r}
  (C\tau)^r
  \left(
    F_{n+r}^N(a)
    +
    \varepsilon_{N,n+r}
  \right)
  +
  2\binom{n+K}{K+1}(C\tau)^{K+1}.
\end{aligned}
  \label{eq:skew-shifted-iteration}
\end{equation}
Now choose \(\tau>0\) such that \(C\tau<1\), and let \(M\in\mathbb N\) be such that \(T\le M\tau\). We prove by induction on \(p=0,\ldots,M\) that, for every fixed \(n\ge1\),
\[
  \lim_{N\to\infty}
  \sup_{0\le t\le p\tau}
  F_n^N(t)
  =
  0.
\]
For \(p=0\), this follows from the initial chaoticity assumption, since \(F_n^N(0)=\delta_N^n\). Assume that the claim holds for some \(p<M\). Fix \(n\ge1\). Applying \eqref{eq:skew-shifted-iteration} with \(a=p\tau\), we obtain, for every fixed \(K\),
\[
  \limsup_{N\to\infty}
  \sup_{p\tau\le t\le(p+1)\tau}
  F_n^N(t)
  \le
  2\binom{n+K}{K+1}(C\tau)^{K+1},
\]
because the induction hypothesis gives \(F_{n+r}^N(p\tau)\to0\) for each fixed \(r=0,\ldots,K\), and because \(\varepsilon_{N,n+r}\to0\). Letting \(K\to\infty\), and using \(C\tau<1\), gives
\[
  \lim_{N\to\infty}
  \sup_{p\tau\le t\le(p+1)\tau}
  F_n^N(t)
  =
  0.
\]
Together with the induction hypothesis on \([0,p\tau]\), this proves the claim at level \(p+1\). Since \(T\le M\tau\), we conclude that, for every fixed \(n\ge1\),
\[
  \lim_{N\to\infty}
  \sup_{0\le t\le T}
  F_n^N(t)
  =
  0.
\]
This completes the proof.
\end{proof}


\begin{thebibliography}{99}

\bibitem{Amini2025}
H.~Amini, N.~H. Amini, S.~Chalal and G.~Guo.
\newblock Graphon quantum filtering systems.
\newblock {\em Preprint}, arXiv:2506.12249, 2025.

\bibitem{Ammari2009}
Z.~Ammari and F.~Nier.
\newblock Mean field limit for bosons and propagation of Wigner measures.
\newblock {\em J. Math. Phys.}, 50 (2009), 042107.
\newblock doi:10.1063/1.3115046.

\bibitem{bardos2000}
C.~Bardos, F.~Golse and N.~J. Mauser.
\newblock Weak coupling limit of the \(N\)-particle Schr\"odinger equation.
\newblock {\em Methods Appl. Anal.}, 7 (2000), 275--293.
\newblock doi:10.4310/MAA.2000.v7.n2.a2.

\bibitem{bardos2002}
C.~Bardos, L.~Erd\H{o}s, F.~Golse, N.~J. Mauser and H.-T. Yau.
\newblock Derivation of the Schr\"odinger--Poisson equation from the quantum
\(N\)-body problem.
\newblock {\em C. R. Acad. Sci. Paris, S\'er. I Math.}, 334 (2002), 515--520.
\newblock doi:10.1016/S1631-073X(02)02253-7.

\bibitem{bouten2007}
L.~Bouten, R.~van Handel and M.~R. James.
\newblock An introduction to quantum filtering.
\newblock {\em SIAM J. Control Optim.}, 46 (2007), 2199--2241.
\newblock doi:10.1137/060651239.

\bibitem{bouard2026}
A.~de Bouard, G.~Guo and T.~H\'erouard.
\newblock Infinite dimensional mean-field Belavkin equation: well-posedness
and derivation.
\newblock {\em Ann. Henri Poincar\'e}, 2026.
\newblock doi:10.1007/s00023-026-01700-4.

\bibitem{frohlich2007}
J.~Fr\"ohlich, S.~Graffi and S.~Schwarz.
\newblock Mean-field and classical limit of many-body Schr\"odinger dynamics
for bosons.
\newblock {\em Comm. Math. Phys.}, 271 (2007), 681--697.
\newblock doi:10.1007/s00220-007-0207-5.

\bibitem{golsepaul2017}
F.~Golse and T.~Paul.
\newblock The Schr\"odinger equation in the mean-field and semiclassical
regime.
\newblock {\em Arch. Ration. Mech. Anal.}, 223 (2017), 57--94.
\newblock doi:10.1007/s00205-016-1031-x.

\bibitem{golsepaul2019}
F.~Golse and T.~Paul.
\newblock Empirical measures and quantum mechanics: applications to the
mean-field limit.
\newblock {\em Comm. Math. Phys.}, 369 (2019), 1021--1053.
\newblock doi:10.1007/s00220-019-03357-z.

\bibitem{Guo2026}
G.~Guo, H.~Liang and Z.~Wang.
\newblock Quantum relative entropy and the mean-field limit.
\newblock {\em Preprint}, arXiv:2605.08652, 2026.

\bibitem{knowlespickl2010}
A.~Knowles and P.~Pickl.
\newblock Mean-field dynamics: singular potentials and rate of convergence.
\newblock {\em Comm. Math. Phys.}, 298 (2010), 101--138.
\newblock doi:10.1007/s00220-010-1010-2.

\bibitem{kolokoltsov2021}
V.~N. Kolokoltsov.
\newblock The law of large numbers for quantum stochastic filtering and
control of many-particle systems.
\newblock {\em Theoret. Math. Phys.}, 208 (2021), 937--957.
\newblock doi:10.1134/S0040577921070084.

\bibitem{kolokoltsov2022}
V.~N. Kolokoltsov.
\newblock Quantum mean-field games.
\newblock {\em Ann. Appl. Probab.}, 32 (2022), 2254--2288.
\newblock doi:10.1214/21-AAP1733.

\bibitem{kolokoltsov2025a}
V.~N. Kolokoltsov.
\newblock On quantum stochastic master equations.
\newblock {\em Electron. J. Probab.}, 30 (2025), article no.~57, 1--21.
\newblock doi:10.1214/25-EJP1322.

\bibitem{kolokoltsov2025b}
V.~N. Kolokoltsov.
\newblock On the mathematical theory of quantum stochastic filtering equations
for mixed states.
\newblock {\em Russ. J. Math. Phys.}, 32 (2025), 510--529.
\newblock doi:10.1134/S1061920825600825.

\bibitem{kolokoltsov2026survey}
V.~N. Kolokoltsov.
\newblock Quantum filtering and propagation of chaos for open quantum systems,
with applications to quantum feedback control and quantum mean-field games.
\newblock {\em Preprint},  2026.

\bibitem{mirrahimivanhandel2007}
M.~Mirrahimi and R.~van Handel.
\newblock Stabilizing feedback controls for quantum systems.
\newblock {\em SIAM J. Control Optim.}, 46 (2007), 445--467.
\newblock doi:10.1137/050644793.

\bibitem{mora2008}
C.~M. Mora and R.~Rebolledo.
\newblock Basic properties of nonlinear stochastic Schr\"odinger equations
driven by Brownian motions.
\newblock {\em Ann. Appl. Probab.}, 18 (2008), 591--619.
\newblock doi:10.1214/105051607000000311.

\bibitem{nielsenchuang2010}
M.~A. Nielsen and I.~L. Chuang.
\newblock {\em Quantum Computation and Quantum Information}.
\newblock 10th anniversary ed., Cambridge University Press, 2010.

\bibitem{pellegrini2008}
C.~Pellegrini.
\newblock Existence, uniqueness and approximation of a stochastic
Schr\"odinger equation: the diffusive case.
\newblock {\em Ann. Probab.}, 36 (2008), 2332--2353.
\newblock doi:10.1214/08-AOP391.

\bibitem{pickl2011}
P.~Pickl.
\newblock A simple derivation of mean-field limits for quantum systems.
\newblock {\em Lett. Math. Phys.}, 97 (2011), 151--164.
\newblock doi:10.1007/s11005-011-0470-4.

\bibitem{spohn1980}
H.~Spohn.
\newblock Kinetic equations from Hamiltonian dynamics: Markovian limits.
\newblock {\em Rev. Modern Phys.}, 52 (1980), 569--615.
\newblock doi:10.1103/RevModPhys.52.569.

\bibitem{watrous2018}
J.~Watrous.
\newblock {\em The Theory of Quantum Information}.
\newblock Cambridge University Press, 2018.

\end{thebibliography}

\end{document}
